% =====================================================================
%  ingiliz_empirizmi_v3.tex
%  Revised manuscript after INTEGRATION_NOTE (2026-08-17)
%
%  The Limits of Formal Methods: Notes on the Formal Representation of
%  Stoic Katalepsis and Hume's Critique of Causation
%
%  Changes vs. v2 (PDF ingiliz_empirizmi_v2.pdf):
%    - §2 is SUBSTITUTED by the canonical core_section.tex (same folder;
%      \input below). The old single-sorted language
%      L = {HCB, CC, J, Kat, Gr, Stoic} is superseded by the many-sorted
%      L_0 / L^+; the J(x)->¬HCB(x) axiom is removed (H-I is a candidate,
%      not a default); (Exc1) is derived using M_0a /\ M_0b (Prop. 1).
%    - Abstract and Özet rewritten in implicit-definability language.
%    - §3 gains the four-level Humean analysis (table).
%    - §6 / §7 updated to the narrower claim (L_0, admissible class K).
%    - §8 Appendix, Provenance Table and References updated.
%    - V5 (2026-08-17): §2 gains encoding-sensitivity (L0^A / L0^B),
%      the E0/E1/E2 minimal-enlargement benchmark, and the Gate 1.5
%      non-triviality check; a new §Objections and Replies (with the
%      negative-result matrix) is inserted before the Conclusion; an
%      Open Science Statement is added before the References.
%    - §1, §4, §5 unchanged in substance (cross-references aligned).
%
%  REQUIRES (host preamble): amsthm + proposition/lemma/definition/remark
%  environments, defined below. Compile with:
%      pdflatex ingiliz_empirizmi_v3.tex   (run twice for refs)
% =====================================================================
\documentclass[11pt]{article}
\usepackage[T1]{fontenc}
\usepackage[utf8]{inputenc}
\usepackage{amsmath,amssymb,amsthm}
\usepackage[margin=1in]{geometry}
\usepackage{array,longtable,booktabs}
\usepackage[hidelinks,colorlinks=true,linkcolor=blue,citecolor=blue,urlcolor=blue]{hyperref}
\usepackage{graphicx}

% ---- REQUIRES: theorem-like environments used by core_section.tex ----
\theoremstyle{plain}
\newtheorem{proposition}{Proposition}[section]
\newtheorem{lemma}[proposition]{Lemma}
\theoremstyle{definition}
\newtheorem{definition}[proposition]{Definition}
\theoremstyle{remark}
\newtheorem{remark}[proposition]{Remark}
% amsthm already provides theorem, proof.

\title{The Limits of Formal Methods: Notes on the Formal Representation
of Stoic Katalepsis and Hume's Critique of Causation\\[4pt]
\large Bi\c{c}imsel Ara\c{c}lar{\i}n S{\i}n{\i}r{\i}: Stoac{\i}
Katalepsis ile Humecu Nedensellik Ele\c{s}tirisinin Bi\c{c}imsel
Temsili \"Uzerine Notlar}
\author{Corresponding author}
\date{2026}

\begin{document}
\maketitle

\begin{abstract}
This paper is a methodological note, not a historical thesis. It
examines a narrow question: what can and cannot be represented when
Stoic katalepsis --- the ``cognition'' (Long \& Sedley 1987, 40) of a
kataleptic impression --- and Hume's account of causation are encoded
in first-order logic (FOL)? We formalize the two positions at the
level of their justification conditions and argue that a purely
extensional FOL encoding can represent the material content of both
positions but cannot fix their explanatory and modal residue:
(i)~Hume's claim that a causal belief is unjustified \emph{because} it
is produced by custom has the form of an explanatory (grounding)
claim; we show that the grounding atom
$G(\mathrm{custFact}(b),\mathrm{nonjustFact}(b,c))$ is \emph{not
implicitly definable} in the extensional signature $L_0$ over the
admissible class of reconstructions $\mathcal K$ --- a singular
encoding limit, not a universal claim about the expressive power of
first-order logic;
(ii)~the ancient definition of the kataleptic impression contains a
literal modal clause --- the impression is ``of such a kind as could
not arise from what is not'' (Diogenes Laertius 7.46; Sextus
Empiricus, \emph{Adversus Mathematicos} 7.248, 7.402) --- which any
FOL encoding must de-modalize. The two residues are not symmetric: the
Stoic residue is modal and is recovered by adding a modal operator
$\Box_s$, while the Humean residue is explanatory and requires a
hyperintensional connective (grounding); the negative result is that
neither is fixed by the $L_0$-reduct of the extensional encoding. We
show this by a model-pair argument over $\mathcal K$: two enriched
structures with identical $L_0$-reducts realize different grounding
readings, and no first-order theory in the unexpanded signature can
separate them. We do not claim that formalization ``necessarily''
fails; we claim only that this encoding, in the explicitly specified
extensional signature $L_0$ and admissible class $\mathcal K$, at this
level of granularity, is semantically --- not syntactically ---
underdetermined with respect to the intended normative reading, and
that representing the distinction requires additional structure --- a
justification operator (justification logic: Artemov \& Fitting 2019),
a grounding connective (Fine 2012), or a modal operator beyond $L_0$.
We then correct two historiographical errors that frequently accompany
such formalizations and add documented channels: (i)~the claim that
Hume rejects the Principle of Sufficient Reason must be qualified
--- Hume denies that the causal maxim is rationally demonstrable
(T 1.3.3; SBN 78--82) and denies any ``absolute or metaphysical
necessity'' in causal connection (T 1.3.14.35; SBN 172);
(ii)~the transmission of Stoic epistemology to early modern Europe ran
through documented printed texts --- Sextus Empiricus (Latin
\emph{Hypotyposes}, tr.\ Henri Estienne, 1562; Latin \emph{Adversus
Mathematicos}, tr.\ Gentian Hervet, 1569; Greek \emph{Opera}, 1621),
Diogenes Laertius (Traversari's translation completed 1433; Latin
\emph{editio princeps} Rome 1472; Venice 1475; Greek \emph{editio
princeps} Basel 1533), and Cicero's \emph{Academica} (\emph{editio
princeps} Rome 1471; the fist analogy at \emph{Lucullus} 145 = Long \&
Sedley 41B), supplemented by the Neostoic reception (Lipsius 1604)
--- not through a speculative medieval chain. The paper concludes that
formal tools are useful precisely when their limits are stated, that
the history of philosophy requires documentary discipline that formal
notation cannot supply, and that this stance has systematic analogues
in non-Western traditions of textual and logical rigor (the Buddhist
\emph{catu\c{s}ko\d{t}i}; the Xunzian theory of the rectification of
names; the Qing \emph{kaozheng} tradition), which we invoke as
analogies, not as transmission claims.
\end{abstract}

\noindent\textbf{Keywords:} Stoicism; katalepsis; phantasia kataleptike;
Hume; causation; Principle of Sufficient Reason; first-order logic;
justification logic; grounding; implicit definability; transmission of
ancient philosophy; catu\c{s}ko\d{t}i; zhengming; kaozheng

\vspace{1em}

\begin{center}\textbf{\"{O}zet}\end{center}
Bu makale tarihsel bir tez de\u{g}il, \textbf{metodolojik bir nottur}.
Dar bir soruyu inceliyoruz: Stoac{\i} \emph{katalepsis} (kavray{\i}c{\i}
izlenimin ``kavranmas{\i}''; Long \& Sedley 1987'nin \c{c}evirisiyle
``cognition'') ile Hume'un nedensellik a\c{c}{\i}klamas{\i} birinci
derece mant{\i}kta (FOL) kodland{\i}\u{g}{\i}nda ne temsil edilebilir,
ne temsil edilemez? \.{I}ki konumu gerek\c{c}elendirme ko\c{s}ullar{\i}
d\"uzeyinde bi\c{c}imselle\c{s}tiriyor ve tamamen ekstansiyonel bir FOL
kodlamas{\i}n{\i}n her iki konumun \textbf{maddi i\c{c}eri\u{g}ini}
temsil edebildi\u{g}ini, ancak \textbf{a\c{c}{\i}klay{\i}c{\i} ve modal
kal{\i}nt{\i}y{\i}} sabitleyemedi\u{g}ini g\"osteriyoruz:
(i)~Hume'un ``nedensel inan\c{c}, al{\i}\c{s}kanl{\i}k taraf{\i}ndan
\"uretildi\u{g}i \emph{i\c{c}in} gerek\c{c}elendirilmemi\c{s}tir''
iddias{\i} a\c{c}{\i}klay{\i}c{\i} (grounding) bi\c{c}imindedir;
$G(\mathrm{custFact}(b),\mathrm{nonjustFact}(b,c))$ grounding atomunun
$L_0$ imzas{\i} ve $\mathcal K$ kabul edilebilir yap{\i} s{\i}n{\i}f{\i}
\"uzerinde \textbf{"ort\"uk olarak tan{\i}mlanabilir (implicitly
definable) olmad{\i}\u{g}{\i}n{\i}} g\"osteriyoruz --- birinci derece
mant{\i}\u{g}{\i}n ifade g\"uc\"u hakk{\i}nda evrensel de\u{g}il, tekil
bir kodlama s{\i}n{\i}r{\i};
(ii)~kavray{\i}c{\i} izlenimin antik tan{\i}m{\i} lafz{\i}
(s\"ozc\"u\u{g}\"u s\"ozc\"u\u{g}\"une, ger\c{c}ek anlamda) bir modal
kloz i\c{c}erir --- izlenim ``olmayandan do\u{g}amayacak
t\"urden''dir (Diogenes Laertius 7.46; Sextus Empiricus,
\emph{Adversus Mathematicos} 7.248, 7.402) --- ve her FOL kodlamas{\i}
bu klozu de-modalize etmek zorundad{\i}r. \.{I}ki kal{\i}nt{\i}
simetrik de\u{g}ildir: Stoac{\i} kal{\i}nt{\i} modaldir ve bir modal
operat\"or ($\Box_s$) eklenerek geri kazan{\i}l{\i}r; Humecu kal{\i}nt{\i}
a\c{c}{\i}klay{\i}c{\i}d{\i}r ve hiperintansiyonel bir ba\u{g}la\c{c}
(grounding) gerektirir; olumsuz sonu\c{c}, ikisinin de ekstansiyonel
kodlaman{\i}n $L_0$-reduktu taraf{\i}ndan sabitlenmemi\c{s} olmas{\i}d{\i}r.
Bunu $\mathcal K$ \"uzerinde bir model-\c{c}ifti arg\"uman{\i}y la
g\"osteriyoruz: \"ozde\c{s} $L_0$-reduktlar{\i}na sahip iki geni\c{s}letilmi\c{s}
yap{\i} farkl{\i} grounding okumalar{\i} ger\c{c}ekle\c{s}tirir ve
geni\c{s}letilmemi\c{s} $L_0$ imzas{\i}nda hi\c{c}bir birinci derece
teori bunlar{\i} ay{\i}rt edemez. Formalizasyonun ``zorunlu olarak''
ba\c{s}ar{\i}s{\i}z oldu\u{g}unu iddia etmiyoruz; yaln{\i}zca bu
kodlaman{\i}n, a\c{c}{\i}k\c{c}a belirlenmi\c{s} $L_0$ imzas{\i} ve
$\mathcal K$ yap{\i} s{\i}n{\i}f{\i} alt{\i}nda, bu ayr{\i}nt{\i}
d\"uzeyinde, ama\c{c}lanan normatif okumaya g\"ore \textbf{semantik
olarak} (s\"ozdizimsel de\u{g}il) eksik belirlenmi\c{s} oldu\u{g}unu
ve ayr{\i}m{\i}n temsili i\c{c}in $L_0$ \"otesinde ek yap{\i} --- bir
gerek\c{c}e operat\"or\"u (gerek\c{c}e mant{\i}\u{g}{\i}: Artemov \&
Fitting 2019), bir grounding ba\u{g}la\c{c}{\i} (Fine 2012) veya bir
modal operat\"or --- gerekti\u{g}ini s\"oyl\"uyoruz. Ard{\i}ndan bu
t\"ur bi\c{c}imselle\c{s}tirmelere s{\i}kl{\i}kla e\c{s}lik eden iki
tarihsel hatay{\i} d\"uzeltiyor ve belgelenmi\c{s} kanallar{\i}
ekliyoruz: (i)~Hume'un Yeter Sebep \.{I}lkesi'ni reddetti\u{g}i iddias{\i}
nitelendirilmelidir --- Hume, nedensel \"ozdeyi\c{s}in rasyonel olarak
g\"osterilemeyece\u{g}ini (T 1.3.3; SBN 78--82) ve nedensel
ba\u{g}lant{\i}da hi\c{c}bir ``mutlak ya da metafizik zorunluluk''
olmad{\i}\u{g}{\i}n{\i} (T 1.3.14.35; SBN 172) s\"oyler;
(ii)~Stoac{\i} epistemolojinin erken modern Avrupa'ya aktar{\i}m{\i}
belgelenebilir bas{\i}l{\i} metinler \"uzerinden ger\c{c}ekle\c{s}mi\c{s}tir
--- Sextus Empiricus (Latince \emph{Hypotyposes}, \c{c}ev. Henri
Estienne, 1562; Latince \emph{Adversus Mathematicos}, \c{c}ev. Gentian
Hervet, 1569; Yunanca \emph{Opera}, 1621), Diogenes Laertius
(Traversari'nin \c{c}evirisi 1433'te tamamland{\i}; Latince \emph{editio
princeps} Roma 1472; Venedik 1475; Yunanca \emph{editio princeps} Basel
1533) ve Cicero'nun \emph{Academica}'s{\i} (\emph{editio princeps} Roma
1471; yumruk benzetmesi \emph{Lucullus} 145 = Long \& Sedley 41B) ile
Neostoac{\i} al{\i}mlan{\i}\c{s} (Lipsius 1604) --- spek\"ulatif bir
orta\c{c}a\u{g} zinciri \"uzerinden de\u{g}il. Makale, bi\c{c}imsel
ara\c{c}lar{\i}n s{\i}n{\i}rlar{\i} belirtildi\u{g}inde yararl{\i}
oldu\u{g}u, felsefe tarihinin bi\c{c}imsel notasyonun sa\u{g}layamad{\i}\u{g}{\i}
bir belge disiplini gerektirdi\u{g}i ve bu tutumun Bat{\i} d{\i}\c{s}{\i}
geleneklerde de sistematik kar\c{s}{\i}l{\i}klar{\i}n{\i}n bulundu\u{g}u
(Budist \emph{catu\c{s}ko\d{t}i}; Xunzici adlar{\i}n d\"uzeltilmesi
teorisi; Qing \emph{kaozheng} gelene\u{g}i) sonucuna var{\i}r --- bu
kar\c{s}{\i}l{\i}klar aktar{\i}m iddias{\i} de\u{g}il, analoji olarak
an{\i}l{\i}r.

\noindent\textbf{Anahtar Kelimeler:} Stoac{\i}l{\i}k; \emph{katalepsis};
\emph{phantasia kataleptike}; Hume; nedensellik; Yeter Sebep \.{I}lkesi;
birinci derece mant{\i}k; gerek\c{c}e mant{\i}\u{g}{\i}; grounding;
\"ort\"uk tan{\i}mlanabilirlik; antik felsefenin aktar{\i}m{\i};
catu\c{s}ko\d{t}i; zhengming; kaozheng

\tableofcontents

% =====================================================================
\section{Introduction: Scope and Limits}
\label{sec:intro}

The purpose of this note can be summarized in three points.

\begin{enumerate}
\item To encode the justification conditions of Stoic katalepsis and
Hume's account of causation at first-order level, and to state
precisely which distinction this encoding misses --- not the material
content of the two positions, but their explanatory and modal residue
(\S\ref{sec:formalization}).
\item To correct two historiographical errors that frequently accompany
such formalizations --- the unqualified claim about Hume and the
Principle of Sufficient Reason (\S\ref{sec:hume-psr}), and the
speculative reconstruction of the transmission of Stoic epistemology
(\S\ref{sec:transmission}) --- and to add the documented channels and
the relevant secondary literature.
\item To avoid universal claims such as ``formalization necessarily
fails'' and to state only singular, verifiable limits --- of this
encoding, at this level of granularity --- together with a statement of
what a richer logical language would add (\S\ref{sec:richer-languages}).
\end{enumerate}

A brief comparative note (\S\ref{sec:comparative}) places the theme of
the limits of formal schemes in a wider context: the theorization of
such limits is not peculiar to the Western analytic tradition. The note
ends with an explicit statement of its own limitations
(\S\ref{sec:limits}) and a provenance table recording the source of
every substantive claim.

Throughout, we observe the following discipline: no citation is
included without a locatable primary text or a verifiable secondary
source; where a claim rests on secondary literature alone, this is
stated; and no claim of the form ``there is no evidence for $X$'' is
allowed to function as ``there is evidence that not-$X$'' (see
\S\ref{sec:limits}).

% =====================================================================
\section{Formalization: What Can and Cannot Be Encoded}
\label{sec:formalization}

% --- SUBSTITUTED: the old single-sorted language
%     L = {HCB, CC, J, Kat, Gr, Stoic} is replaced by the canonical
%     many-sorted L_0 / L^+ core section. The J(x)->¬HCB(x) axiom is
%     gone (H-I is a candidate, not a default); (Exc1) is derived with
%     M_0a /\ M_0b (Proposition 1); the model-pair argument is rebuilt
%     as an implicit-definability failure over the admissible class K.
% Core formal section for
% "What an Extensional First-Order Formalization Leaves Underdetermined:
%  Stoic Katalepsis and Humean Custom"
%
% Revision: 2026-08-17
%   - P0: thesis narrowed to "the target relation is not determined by the
%         L_0-reduct over the admissible class K";
%   - P0: Proposition 1 (formerly "Theorem 1") explicitly conditional;
%   - P0: Bridge-collapse proposition and characterization added;
%   - P0: Proposition 2 rebuilt as an implicit-definability failure,
%         with definability, stability and Obtains sub-sort;
%   - P0: Hume reading labelled "Strong exclusion reconstruction";
%   - P0: Stoic formula relabelled "minimal dependency skeleton";
%   - P1: language-level discipline (FOL, metalanguage, modal, JL);
%   - P1: separate Cont / Fact sorts with Obtains monad;
%   - P1: modal operator \Box_s indexed explicitly;
%   - P1: provenance table updated; comparative section deferred to appendix.
%
% Intended for integration into the main manuscript as \S2 of §2.

\subsection{Interpretive and Formal Preconditions}
\label{sec:formal-preconditions}

This is a methodological note, not a historical thesis. Its aim is to
state precisely what an extensional first-order encoding at a fixed
level of granularity can and cannot represent, and to keep the resulting
claims conditional on the encoding choices that produced them. We do
not attribute a contemporary grounding relation, a fully articulated
theory of epistemic justification, or a complete process-reliabilistic
framework to the historical authors. Where we use such vocabulary, it
is a contemporary regimentation of one strong interpretive option; an
alternative reading is recorded alongside it.

The formal reconstruction is deliberately narrower than a reconstruction
of either Stoic epistemology or Hume's epistemology as a whole. Its
purpose is to isolate the minimum content-bearing structure required for
the representational argument. Accordingly, we bracket the full Stoic
ontology of \emph{lekta} and retain only the structure needed to
distinguish impressions, represented contents, cognitive episodes,
belief-production, and justification. In particular, the predicates
introduced below do not exhaust, nor do they uniquely capture, the
historical vocabularies.

\begin{quote}
For purposes of reconstruction, we bracket the full Stoic ontology of
\emph{lekta} and Hume's full apparatus of passions, sympathy and
moral psychology, and represent only the minimal content-bearing
structure required for the formal result. Where an interpretive
choice is required, the choice is stated explicitly and an
alternative reading is acknowledged.
\end{quote}

\subsection{The Minimal Extensional Language $L_0$}
\label{sec:l0}

Let $L_0$ be a many-sorted first-order language with the following
core sorts:
\[
I \quad\text{(impressions)},\qquad
\mathit{Cont} \quad\text{(ordinary representational contents)},\qquad
B \quad\text{(cognitive episodes)}.
\]
A subsort $\mathit{Fact} \subseteq \mathit{Cont}$ is introduced only in
the enriched language $L^+$ of \S\ref{sec:enrichment} to host
reified grounding \emph{relata}; it is not part of $L_0$ proper.
An additional sort $E$ for explicit justification terms is introduced
only in the justification-logic enrichment of
\S\ref{sec:richer-languages} and is not part of the core extensional
language.

The non-logical vocabulary of $L_0$ is:

\begin{align*}
&\operatorname{Kat}(i)
&&\text{: $i\in I$ is kataleptic;} \\
&\operatorname{Rep}(i,c)
&&\text{: impression $i$ represents content $c\in\mathit{Cont}$;} \\
&\operatorname{Grasp}(b,i)
&&\text{: cognitive episode $b$ grasps impression $i$;} \\
&\operatorname{Assent}(b,c)
&&\text{: $b$ assents to content $c$;} \\
&\operatorname{Bel}(b,c)
&&\text{: $b$ is a belief state with content $c$ (whether or not
causally formed);} \\
&\operatorname{Causal}(b,c)
&&\text{: $b$ is a causal belief with content $c$;} \\
&\operatorname{Custom}(b)
&&\text{: $b$ is produced by custom;} \\
&\operatorname{Just}(b,c)
&&\text{: the relation of $b$ to $c$ is rationally
justified;}\footnotemark \\
&\operatorname{StoicEp}(b)
&&\text{: $b$ is a cognitive episode to which the reconstructed
Stoic condition applies.}
\end{align*}
\footnotetext{$\operatorname{Just}(b,c)$ is binary rather than unary. The
claim that a cognitive episode is ``justified'' \emph{simpliciter} is
distinct from the claim that its assent to a particular content is
justified. The choice of $b$ as ``episode'' (rather than
``assent-token'' or ``agent-relative state'') is itself a hermeneutic
decision, recorded here so that no later inference silently presupposes
a particular reading of Hume's ``belief''.}

This separation is intentional. In particular,
$\operatorname{Grasp}(b,i)$ is not conflated with representation
$\operatorname{Rep}(i,c)$ or assent $\operatorname{Assent}(b,c)$:
$\operatorname{Grasp}$ operates in the impression-impression space,
$\operatorname{Assent}$ in the content space, and $\operatorname{Bel}$
in the cognitive-state space. The distinction prevents later
inferences from silently collapsing what an extensional encoding
should keep separate.

\subsubsection{The Stoic Minimal Dependency Skeleton}
\label{sec:stoic-skeleton}

The Stoic formula is presented here as a \emph{minimal dependency
skeleton}: a relational pattern extracted from the katalepsis
discussion that is sufficient to register the relevant dependency
between kataleptic impression, representation, grasping and assent,
without claiming to exhaust Stoic epistemology. Modal and doxastic
content is recorded separately (cf.\ \S\ref{sec:stoic-modal-note}).

\[
\forall b\,
\Bigl(
\operatorname{StoicEp}(b)
\rightarrow
\exists i\,\exists c\,
\bigl[
\operatorname{Kat}(i)\land
\operatorname{Rep}(i,c)\land
\operatorname{Grasp}(b,i)\land
\operatorname{Assent}(b,c)
\bigr]
\Bigr).
\tag{S}
\]

Formula (S) is a reconstruction constraint, not a claim that the
complete Stoic epistemology reduces to this formula; in particular,
(S) does not by itself guarantee that $b$ is a belief state about
$c$, that the \emph{content} $c$ figured in both $\operatorname{Rep}$
and $\operatorname{Assent}$ is the same content, or that
$\operatorname{Kat}(i)$ is connected to a truth or ``could not arise
from what is not'' clause. The first of these is a deliberate choice:
the link from the minimal skeleton to genuine belief is the place
where the modern interpreter is owed a separate, content-level
argument, and we do not pre-empt that argument here.

\subsubsection{Humean Mechanism and the Exclusion Readings}
\label{sec:hume-readings}

The core formalization records two mechanism-level constraints:

\[
\forall b\,\forall c\,
\bigl(\operatorname{Causal}(b,c)\rightarrow \operatorname{Custom}(b)\bigr),
\tag{M$_{0a}$}
\]
\[
\forall b\,\forall c\,
\bigl(\operatorname{Causal}(b,c)\rightarrow \operatorname{Bel}(b,c)\bigr).
\tag{M$_{0b}$}
\]

No axiom is imposed that identifies custom-production with the
\emph{absence} of justification. Instead, three readings are
distinguished, each as a candidate reconstruction.

\paragraph{Strong exclusion reconstruction (H-I).}
This is one contemporary regimentation of the strongest reading of
Hume's claim that custom-production \emph{excludes} rational
justification:
\[
\forall b\,\forall c\,
\bigl[
\operatorname{Custom}(b)\land
\operatorname{Bel}(b,c)
\rightarrow
\neg\operatorname{Just}(b,c)
\bigr].
\tag{H-I}
\]
H-I is a candidate, not a presupposition: a contemporary interpreter
who takes Hume's talk of ``custom'' to carry this strong normative
weight will write (H-I) into her reconstruction; one who does not, will
not. The choice is hermeneutic, not formal.

\paragraph{Naturalistic-compatible reading (H-II).}
The weaker \emph{naturalistic-compatible} reading imposes no such
axiom. On H-II, $\operatorname{Custom}(b)$ records a
belief-forming mechanism while the justificatory status of the
relation between $b$ and $c$ remains separately determined. The
distinction prevents the illicit inference from psychological genesis
to epistemic invalidity. We treat (H-II) as the default for the rest
of this section, in line with the methodological note's commitment
not to over-read Hume.

\paragraph{Direct-attachment reading ($T_1$).}
For comparison, define
\[
\forall b\,\forall c\,
\bigl[
\operatorname{Causal}(b,c)
\rightarrow
\neg\operatorname{Just}(b,c)
\bigr].
\tag{T$_1$}
\]
$T_1$ attaches the exclusion directly to causal beliefs; it is
logically stronger than H-I only when a further bridge axiom is
present (cf.\ Proposition~\ref{prop:bridge-collapse}).

The next result concerns the extensional relationship between
these formulations.

\subsection{Proposition~1: Strength Relation Between Two Reconstructions}
\label{sec:prop1}

\begin{proposition}[Strength relation between two proposed reconstructions]
\label{prop:strength}
Let $M_0$ be the conjunction of (M$_{0a}$) and (M$_{0b}$), let
$T_1$ be the formula above, and let $T_2$ denote (H-I). Then
\[
T_2\land M_0 \models T_1,
\]
whereas
\[
T_1\land M_0 \not\models T_2.
\]
\end{proposition}

\begin{proof}
For the first entailment, let $\mathcal M\models T_2\land M_0$ and
suppose $\mathcal M\models\operatorname{Causal}(b,c)$. By
(M$_{0a}$), $\mathcal M\models\operatorname{Custom}(b)$, and by
(M$_{0b}$), $\mathcal M\models\operatorname{Bel}(b,c)$. Hence the
antecedent of $T_2$ is true, and since $\mathcal M\models T_2$,
$\mathcal M\models\neg\operatorname{Just}(b,c)$. Therefore
$\mathcal M\models T_1$. Since $b,c$ were arbitrary,
$T_2\land M_0\models T_1$.

For the converse, consider a structure with one cognitive episode
$b_1$ and one content $c_1$ such that
$\operatorname{Causal}(b_1,c_1)=\mathbf{false}$,
$\operatorname{Custom}(b_1)=\mathbf{true}$,
$\operatorname{Bel}(b_1,c_1)=\mathbf{true}$,
$\operatorname{Just}(b_1,c_1)=\mathbf{true}$.
Then (M$_{0a}$) and (M$_{0b}$) are satisfied vacuously, $T_1$ is
satisfied vacuously, but $T_2$ is false (its antecedent is true and
its consequent is false). Hence $T_1\land M_0\not\models T_2$.
\end{proof}

Proposition~1 is purely comparative and conditional. It establishes
neither that Hume endorsed $T_2$, nor that the causal--custom relation
is correctly represented by $M_0$, nor that $T_1$ exhausts
Humean scepticism. The result is about two proposed reconstructions,
relative to a fixed signature, on a fixed class of models.

\subsubsection{Bridge collapse and characterization}
\label{sec:bridge-collapse}

A reader may object that the difference between $T_1$ and $T_2$ is
trivial once a natural interpretive decision is made. The
\emph{bridge axiom}
\[
\forall b\,\forall c\,
\bigl[
(\operatorname{Custom}(b)\land\operatorname{Bel}(b,c))
\rightarrow
\operatorname{Causal}(b,c)
\bigr]
\tag{B$_0$}
\]
encodes the obvious Humean gloss that every custom-produced belief
\emph{is} a causal belief (EHU 5.1; T 1.3.14.20). Under $B_0$, the two
readings collapse.

\begin{proposition}[Bridge collapse]
\label{prop:bridge-collapse}
Over $M_0$,
\[
T_1\land B_0 \models T_2.
\]
Given Proposition~\ref{prop:strength}, $T_1$ and $T_2$ are therefore
extensionally \emph{equivalent} over $M_0\land B_0$. Moreover, at a
fixed pair $(b,c)$ the two readings differ---i.e.\ $T_1$ holds at the
pair while $T_2$ fails there---if and only if
\[
\operatorname{Custom}(b)\land
\operatorname{Bel}(b,c)\land
\operatorname{Just}(b,c)\land
\neg\operatorname{Causal}(b,c).
\tag{$\star$}
\]
Consequently, over $M_0$, a model separates $T_1$ from $T_2$ (i.e.\
satisfies $T_1\land M_0\land\neg T_2$) if and only if it satisfies
$T_1\land M_0$ together with $\exists b\,\exists c\,(\star)$.
\end{proposition}

\begin{proof}
For the entailment, suppose $\mathcal M\models T_1\land M_0\land B_0$
and $\mathcal M\models\operatorname{Custom}(b)\land\operatorname{Bel}(b,c)$.
By $B_0$, $\mathcal M\models\operatorname{Causal}(b,c)$, and by $T_1$,
$\mathcal M\models\neg\operatorname{Just}(b,c)$. Hence $T_2$ holds.
For the characterization: at a fixed pair $(b,c)$, $T_1$ holds and
$T_2$ fails exactly when $(\star)$ holds there. Reading $(\star)$
existentially, $T_1\land M_0\land\neg T_2$ entails $(\star)$:
$\neg T_2$ yields a pair with $\operatorname{Custom}\land
\operatorname{Bel}\land\operatorname{Just}$, on which $T_1$ forces
$\neg\operatorname{Causal}$. Conversely, $T_1\land M_0$ together with
$(\star)$ entails $\neg T_2$ ($(\star)$ is a witness) and hence entails
$T_1\land M_0\land\neg T_2$; the two directions give the equivalence.
Note that $(\star)$ alone does not entail $T_1$: a model may satisfy
$(\star)$ at one pair while violating $T_1$ at another.
\end{proof}

The computational verification in \S\ref{sec:verification} confirms
both that $T_1\land M_0\land B_0\models T_2$ (zero falsifying
interpretations in the finite Boolean fragment) and that $(\star)$ is
the exact condition of separation at a single pair (globally,
separation is $T_1\land M_0$ together with $\exists b\,\exists c\,
(\star)$, as established above). The choice between $T_1$ and $T_2$
is therefore not settled by the logical facts alone; it is settled
by whether one accepts $B_0$. $B_0$ is an interpretive decision: it
identifies the predicate \texttt{Custom} with a mechanism
representable in the \texttt{Causal} vocabulary. That identification
is exactly the \emph{annotation gap} at the heart of this section.
The same kind of gap will reappear, in a more pointed form, in the
reification argument of \S\ref{sec:enrichment}.

\subsection{Controlled Reification and the Enriched Language}
\label{sec:enrichment}

The central difficulty arises when the intended Humean claim is not
merely that two facts co-occur but that one fact \emph{explains} or
\emph{grounds} another. Grounding is hyperintensional: it is not a
truth-functional connective, and it cannot be typed as a relation
between formulas in the object language. Accordingly, the
enrichment below uses \emph{controlled reification}: the explanatory
relata are first reified as fact-like objects, and grounding is then
typed as a relation between those objects. This is a
representational device, not a metaphysical commitment to grounding
facts as ordinary propositions of the historical theories.

Let the content sort $\mathit{Cont}$ also host, in the enriched
language, an additional subsort $\mathit{Fact}\subseteq\mathit{Cont}$.
A unary monadic predicate $\operatorname{Obtains}$ on $\mathit{Fact}$
marks the fact-like objects that actually obtain:
\[
\operatorname{Obtains}(f) \quad (f\in\mathit{Fact})
\]
is read ``$f$ obtains''. $\operatorname{Obtains}$ is part of $L^+$
but not of $L_0$.

Introduce two function symbols
\[
\operatorname{custFact}:B\to\mathit{Fact},
\qquad
\operatorname{nonjustFact}:B\times\mathit{Cont}\to\mathit{Fact},
\]
with the intended readings:
\begin{quote}
$\operatorname{custFact}(b)$ is the fact that $b$ is
custom-produced; $\operatorname{nonjustFact}(b,c)$ is the fact that
the relation of $b$ to $c$ lacks justification.
\end{quote}
The reification is disciplined by three bridge axioms, all in $L^+$:
\begin{align*}
\operatorname{Obtains}(\operatorname{custFact}(b))
&\leftrightarrow \operatorname{Custom}(b), \tag{O$_1$} \\
\operatorname{Obtains}(\operatorname{nonjustFact}(b,c))
&\leftrightarrow \neg\operatorname{Just}(b,c), \tag{O$_2$} \\
G(x,y) &\rightarrow \operatorname{Obtains}(x)\land\operatorname{Obtains}(y),
\tag{O$_3$}
\end{align*}
where $G$ is a new binary relation
\[
G\subseteq \mathit{Fact}\times\mathit{Fact},
\]
read as ``grounds''. The target statement is therefore the
well-typed atomic formula
\[
G\bigl(\operatorname{custFact}(b),\,
       \operatorname{nonjustFact}(b,c)\bigr).
\tag{G}
\]
(O$_3$) records the standing requirement of contemporary grounding
theories that grounding implies truth (or obtainment) of the
relata. (O$_1$) and (O$_2$) pin down what it is for the reified
fact-objects to obtain; without them, $\operatorname{custFact}$ and
$\operatorname{nonjustFact}$ would be ``floating'' objects that
exist independently of whether the base condition holds, which is
ontologically profligate. The proof of Proposition~\ref{prop:underdet}
does not depend on (O$_1$)---(O$_3$); they are stated here so that
the reification device is not a free-floating formalism.

\subsection{The Class of Admissible Reconstructions $\mathcal K$}
\label{sec:admissible}

For the underdetermination result to be a substantive claim about
\emph{this} encoding rather than a triviality, the class of
admissible reconstructions must be specified. Let
\[
L^+ \;=\; L_0 \,\cup\,
\{\,\operatorname{custFact},\,
  \operatorname{nonjustFact},\,
  \operatorname{Obtains},\,
  G \,\}.
\]
A structure $\mathcal M^+$ for $L^+$ is \emph{admissible} (i.e.\
$\mathcal M^+\in\mathcal K$) iff it satisfies the following
constraints:
\begin{enumerate}
\item $\mathcal M^+$ satisfies the mechanism axioms (M$_{0a}$) and
      (M$_{0b}$) on its $L_0$-reduct;
\item the typing of $\operatorname{custFact}$,
      $\operatorname{nonjustFact}$, $G$ is respected: the function
      symbols land in $\mathit{Fact}$, and $G$ is a binary relation
      on $\mathit{Fact}\times\mathit{Fact}$;
\item the structural axioms (O$_1$)--~(O$_3$) are satisfied;
\item \emph{no further bridge axiom} connects
      $\operatorname{Custom}$, $\operatorname{Just}$, or
      $\operatorname{Bel}$ to the extension of $G$ beyond (O$_1$)
      and (O$_2$). In particular, no axiom of the form
      $\varphi(b,c)\rightarrow G(\operatorname{custFact}(b),
      \operatorname{nonjustFact}(b,c))$, with $\varphi$ an $L_0$
      formula, is imposed.
\end{enumerate}
The fourth condition is the operative one. It is what makes the
underdetermination result interesting: the encoding does not
\emph{by stipulation} settle the question we are trying to put to
it. Other choices of $\mathcal K$ are possible; the result is
relative to the one stated.

\subsection{Proposition~2: Implicit Definability Failure over $\mathcal K$}
\label{sec:prop2}

\begin{lemma}[$L_0$-reduct invariance]
\label{lem:reduct-invariance}
Let $\mathcal M_1^+,\mathcal M_2^+\in\mathcal K$ with identical
$L_0$-reducts:
$\mathcal M_1^+\!\upharpoonright L_0
 =\mathcal M_2^+\!\upharpoonright L_0$.
Then, for every $L_0$-sentence $\varphi$,
$\mathcal M_1^+\models\varphi \Leftrightarrow
 \mathcal M_2^+\models\varphi$.
\end{lemma}

\begin{proof}
Induction on the complexity of $L_0$ formulas. Atomic $L_0$ formulas
have the same truth value in the two structures because their
interpretations are identical on the common $L_0$-reduct. The
Boolean cases follow immediately. The quantified cases follow because
the relevant domains and interpretations of all $L_0$ symbols are the
same in the two structures. $G$, $\operatorname{Obtains}$,
$\operatorname{custFact}$, $\operatorname{nonjustFact}$ are not in
$L_0$ and therefore cannot affect the truth of any $L_0$ sentence.
\end{proof}

\begin{proposition}[Implicit definability failure of the grounding atom]
\label{prop:underdet}
There exist admissible structures $\mathcal M_1^+,\mathcal M_2^+\in\mathcal K$
with identical $L_0$-reducts such that
\[
\mathcal M_1^+\models
G\bigl(\operatorname{custFact}(b),\operatorname{nonjustFact}(b,c)\bigr),
\qquad
\mathcal M_2^+\not\models
G\bigl(\operatorname{custFact}(b),\operatorname{nonjustFact}(b,c)\bigr).
\]
Consequently, the grounding atom (G) is \emph{not implicitly
definable} in $L_0$ over $\mathcal K$.
\end{proposition}

\begin{proof}
Take a fixed admissible $L_0$-structure $\mathcal M$ and construct
\[
\mathcal M_1^+=\langle\mathcal M,\,G_1\rangle,\qquad
\mathcal M_2^+=\langle\mathcal M,\,G_2\rangle,
\]
where $G_1$ contains
$\bigl(\operatorname{custFact}(b),\operatorname{nonjustFact}(b,c)\bigr)$
and $G_2$ does not. Choose $\mathcal M$ so that
$\operatorname{Custom}(b)$ and $\neg\operatorname{Just}(b,c)$ hold; then by
(O$_1$)--(O$_2$) both relata of the target atom obtain, and (O$_3$) holds in
$\mathcal M_1^+$ (and vacuously in $\mathcal M_2^+$). Hence both structures
are admissible, i.e.\ $\mathcal M_1^+,\mathcal M_2^+\in\mathcal K$. By
construction,
$\mathcal M_1^+\!\upharpoonright L_0
 =\mathcal M_2^+\!\upharpoonright L_0$.
By Lemma~\ref{lem:reduct-invariance}, the two structures agree on
every $L_0$-sentence. Nevertheless, they disagree on (G). Hence no
$L_0$-sentence determines whether (G) holds. This is exactly the
failure of implicit definability: the $L_0$-reduct alone does not
fix the extension of $G$ at the target pair.
\end{proof}

\begin{proposition}[Definitional irreducibility of the grounding atom]
\label{prop:definability}
Let $\mathcal M_1^+,\mathcal M_2^+$ be the pair of
Proposition~\ref{prop:underdet}. There is no $L_0$-formula
$\theta(b,c)$ with free variables of sort $B$ and $\mathit{Cont}$ such
that, for both $i=1,2$,
\[
\mathcal M_i^+\models
\forall b\,\forall c\,
\bigl[
\theta(b,c)\leftrightarrow
G(\operatorname{custFact}(b),\operatorname{nonjustFact}(b,c))
\bigr].
\]
\end{proposition}

\begin{proof}
The atom (G) is true in $\mathcal M_1^+$ and false in
$\mathcal M_2^+$, while $\theta$'s interpretation depends only on
the $L_0$-reduct, which is identical in both. Hence $\theta$ has
the same truth value at $(b,c)$ in both structures, and the
biconditional cannot hold in both simultaneously.
\end{proof}

\begin{proposition}[Stability under arbitrary $L_0$-theories]
\label{prop:stability}
Let $T$ be any set of $L_0$-sentences and let
$\mathcal M\models T$. Then the enrichment pair of
Proposition~\ref{prop:underdet} satisfies
$\mathcal M_1^+\models T$ and $\mathcal M_2^+\models T$. Consequently,
no $L_0$-theory can separate the two grounding readings while
preserving the reduct $\mathcal M$; any constraint on $G$ must be
stated in the enriched signature itself.
\end{proposition}

\begin{proof}
The two enrichments share the same $L_0$-reduct, so by
Lemma~\ref{lem:reduct-invariance} they agree on every
$L_0$-sentence; in particular, on every sentence of $T$.
\end{proof}

The three propositions jointly establish the underdetermination
claim at three different levels of strength. The mechanism that
produces them is general: it would also produce them for a relation
$G$ between arbitrary other reified objects, and the present result
is not a \emph{content-specific} limit of the encoding but a
consequence of the simple fact that $G$ is not in $L_0$. The
\emph{interest} of the result is not the mechanism but the
\emph{content} of the symbol on which it acts: $G$ is the one
relation that the contemporary grounding literature standardly
deploys to regiment the kind of explanatory claim at issue in
Hume, and the result shows that regimenting that claim requires
semantic resources beyond the $L_0$-reduct.

\subsubsection{Definability qualification: explicit, implicit, and the Beth anchor}
\label{sec:definability-qual}

The underdetermination result of
\S\ref{sec:enrichment}--\ref{sec:prop2} is best read against the
classical background of definability theory.

\begin{definition}[Explicit definability]
A relation $R$ on a structure $\mathcal M$ is \emph{explicitly
definable in $L_0$} iff there is an $L_0$-formula $\theta(\bar x)$
with $\mathcal M\models\forall\bar x\,(R(\bar x)\leftrightarrow\theta(\bar x))$.
\end{definition}

\begin{definition}[Implicit definability]
A relation $R$ is \emph{implicitly definable in $L_0$ over a class
$\mathcal K$} iff for every $\mathcal M_1,\mathcal M_2\in\mathcal K$,
$\mathcal M_1\!\upharpoonright L_0 = \mathcal M_2\!\upharpoonright L_0$
implies $\mathcal M_1\!\upharpoonright R = \mathcal M_2\!\upharpoonright R$.
\end{definition}

Proposition~\ref{prop:definability} is a failure of explicit
definability; Proposition~\ref{prop:underdet} is a failure of
implicit definability; Proposition~\ref{prop:stability} says the
implicit failure persists under any $L_0$-theory. The
\emph{Beth definability theorem} (Beth 1953) guarantees, under
suitable conditions, that implicit definability entails explicit
definability. The conditions of Beth's theorem are
\emph{not} satisfied in the present setting: the class
$\mathcal K$ is closed under $L_0$-elementary equivalence between
its members but not under the constructions Beth's proof requires
(arbitrary ultrapowers or unions of chains of $L_0$-elementary
substructures). The present argument therefore does \emph{not} rely
on Beth's theorem; it is a direct construction of a pair of
admissible structures with the required properties. The theorem is
mentioned here only to fix vocabulary and to make the
underdetermination claim comparable to the standard model-theoretic
literature.

\subsection{What Richer Languages Add---and Do Not Add}
\label{sec:richer-languages}

The negative result concerns the chosen extensional signature. Any
richer representation must therefore add semantic resources rather
than merely re-interpret the original $L_0$ vocabulary. We record
three such resources; the choice among them is the philosophical
work, and this section is deliberately neutral among them.

\subsubsection{Modal enrichment $\Box_s$}
\label{sec:modal-enrichment}

A modal enrichment introduces a unary operator $\Box_s$ together
with an accessibility relation $R$ on the (extended) domain. The
subscript $s$ is an explicit \emph{index} that fixes the intended
reading: $s$ may be (i) a source-reliability index (the impression's
provenance), (ii) an epistemic accessibility index, or (iii) a
metaphysical necessity index. The choice is interpretive; for
example, the Stoic clause ``could not arise from what is not'' is
naturally read as a \emph{source-reliability} claim and would be
regimented as $\Box_{\mathrm{src}}\,\varphi$, while a metaphysical
necessity reading of the same clause would be $\Box_{\mathrm{met}}\,\varphi$.
Conflating these readings is a common error, and we do not pre-empt
it here. Formulas such as $\Box_s\varphi$ admit a standard
translation into FOL over enriched structures (the standard
translation of modal logic into first-order logic with a binary
accessibility relation); the translation changes the signature, and
even then it captures truth in a model, not the hyperintensional
``because'' (cf.\ \S\ref{sec:enrichment}). Modal enrichment is
\emph{necessary but not sufficient} for representing the
explanatory residue.

\subsubsection{Justification-logic enrichment}
\label{sec:just-enrichment}

A justification-logic enrichment introduces explicit evidence terms
of sort $E$ and expressions of the form
\[
t:\varphi,
\]
read as ``$t$ is a justification for $\varphi$''. The construction
$t:\varphi$ is \emph{not} an $L_0$ formula: it is a separate
syntactic category, governed by a separate proof system, in which
``$t:\varphi$'' is a judgement rather than a proposition. In
particular, the term language has its own formation rules, and the
factivity schema $t:\varphi\rightarrow\varphi$ (or, in weaker
variants, $t:\varphi\rightarrow(\varphi\lor\neg\varphi)$) is an
inference rule of the metalanguage, not a $L_0$ axiom. Conflating
$t:\varphi$ with a $L_0$ predicate ``$t$ justifies $\varphi$'' is a
frequent error and obscures the way in which evidence becomes
representable. With this enrichment, the Humean situation is
represented structurally: for a causal belief $a$, the theory
contains $c:\operatorname{Custom}(a)$ (the custom term attests the
production claim) while, at the meta-level, the theory has no
theorem of the form $r:\operatorname{Just}(a,c)$ for any closed
rational-ground term $r$. The two roles that the extensional
encoding collapses---``what produced the belief'' and ``what
justifies it''---are separated at the level of the term language.
This is the minimal positive moral of our negative result: the
distinction requires a language in which \emph{evidence} is an
explicit first-class object.

\subsubsection{Grounding enrichment}
\label{sec:grounding-enrichment}

A grounding enrichment introduces a partial-grounding connective
$<$ on $\mathit{Fact}\times\mathit{Fact}$ with the standard
regimentations: irreflexivity, asymmetry and transitivity
(Fine 2012; Correia \& Schnieder 2012). Under the reification
device of \S\ref{sec:enrichment}, the Humean exclusion is written
$G(\operatorname{custFact}(b),\operatorname{nonjustFact}(b,c))$ as a
grounding claim, with $G$ playing the role of $<$ in a regimentation
that respects (O$_3$). The direction of determination---the feature
absent from model theory---is the very content of the connective.
Note that grounding is hyperintensional, so the failure of $L_0$ to
fix the extension of $G$ is not specific to FOL's extensionality;
it persists for modal extensions as well. This is why modal
enrichment, while necessary, is not sufficient.

\subsubsection{Disjunctive conclusion}
\label{sec:disjunctive}

We draw no conclusion about which of the three enrichments is
``the'' correct one. The conclusion is disjunctive: \emph{some}
hyperintensional, modal, or explicit-evidence resource is required,
and the $L_0$-reduct is not such a resource at this granularity.
The result is therefore about this encoding, at this level of
granularity, under extensional semantics. It is not a universal
claim about the expressive limits of first-order logic, nor a
denial that any richer first-order language could represent the
relevant distinction.

\subsection{Historical Anchoring: The Four-Level Humean Analysis}
\label{sec:hume-historical}

The Humean readings of \S\ref{sec:hume-readings} are not
historically neutral. To keep the formal result from silently
importing a contested interpretive thesis, we record a four-level
decomposition of the claim and tag each level with its textual
anchoring and its standing in the literature.

\begin{center}
\begin{tabular}{p{0.18\linewidth} p{0.42\linewidth} p{0.30\linewidth}}
\hline
\textbf{Level} & \textbf{Claim} & \textbf{Status} \\
\hline
Psychological-genetic &
Causal expectations are formed by custom/habit, not by demonstrative
reason. &
Directly defensible from Hume's text (T 1.3.14.20, SBN 164--65;
E 7.6, SBN 63; T 1.3.16.9, SBN 178--9). \\
Epistemic &
These beliefs are not demonstratively justified. &
Defensible, but the term ``justification'' must be unpacked
(§\ref{sec:hume-readings}). \\
Normative exclusion &
Custom-production excludes rational justification. &
Contested contemporary reconstruction; alternative readings exist
(H-II, naturalistic; Della Rocca 2010 on PSR as demotion vs\
rejection; Pruss 2006 on Hume as evaporationist). \\
Grounding &
The custom-production relation \emph{grounds} the absence of
justification. &
Contemporary regimentation in the sense of Fine (2012) and
Schnieder (2011); not directly attributable to Hume. \\
\hline
\end{tabular}
\end{center}

The formal apparatus above operates at levels 2--4 only with the
qualifications noted. We do not attribute a contemporary grounding
relation, nor a fully articulated theory of epistemic justification,
to Hume. We use grounding only to regiment one strong contemporary
reconstruction of the explanatory force that an interpreter may
take Hume's appeal to custom to possess; alternative reconstructions
remain live, and the formal apparatus is invariant under the choice
between H-I and H-II modulo the annotations recorded in
\S\ref{sec:prop1}--\S\ref{sec:prop2}.

\subsection{The Stoic Modal Clause: A Note}
\label{sec:stoic-modal-note}

The ancient definition of the kataleptic impression (DL 7.46
[Dorandi 2013]; M 7.248, 7.402 [Bury 1935]) contains a literal modal
clause: the impression is ``of such a kind as could not arise from
what is not''. The minimal dependency skeleton (S) of
\S\ref{sec:stoic-skeleton} \emph{does not} contain this clause.
Encoding it requires either (i) a modal operator $\Box_s$ with an
appropriately indexed reading (§\ref{sec:modal-enrichment}), or
(ii) a separate monadic predicate $\operatorname{Ver}(i)$ on
impressions together with an explicit axiom
$\operatorname{Kat}(i)\rightarrow\operatorname{Ver}(i)$. The choice
between (i) and (ii) is interpretive. We do not pre-empt it: the
present formalization is at the level of the minimal dependency
skeleton, and the modal clause is identified as a separate residue
to be added by whichever enrichment the reader finds defensible.

The separation of (S) from the modal clause is methodologically
important. It is the formal counterpart of the historical point
that katalepsis and \emph{epist\=em\=e} are
\emph{not} identical in Stoic epistemology: katalepsis is the
cognition of a kataleptic impression, while \emph{epist\=em\=e}
is a stable, secure cognitive state distinct from mere cognitive
uptake. The minimal dependency skeleton leaves room for this
distinction; encoding the Stoic position more fully would require
additional predicates and a separate argument that we do not
undertake here.

\subsection{Verification Status and Reproducibility}
\label{sec:verification}

The formal core has been checked at two levels.

\emph{First}, the finite Boolean fragment underlying
Proposition~\ref{prop:strength} and
Proposition~\ref{prop:bridge-collapse} was exhaustively enumerated
for the single-$B$, single-$\mathit{Cont}$ case. Sixteen interpretations were
generated. The consequence $T_2\land M_0\models T_1$ was confirmed
(zero falsifying interpretations); the stated countermodel to
$T_1\land M_0\models T_2$ was verified; the bridge collapse
$T_1\land B_0\models T_2$ was verified (zero falsifying
interpretations in the Boolean fragment); and the characterization
$(\star)$ of the separation between $T_1$ and $T_2$ over $M_0$ was
confirmed. The full table of 16 interpretations is given in
\S\ref{sec:verification-table} below.

\emph{Second}, the model pair required for
Proposition~\ref{prop:underdet} was explicitly constructed with
identical $L_0$ interpretations and distinct $G$ extensions. This
machine check does not replace the general semantic proof: the
latter is supplied by Lemma~\ref{lem:reduct-invariance}. The
computational check is therefore a reproducibility and
error-detection layer, while the lemma-and-proposition argument
remains the mathematical proof presented in this section.

The associated reproducibility materials implement exactly this
finite check and the two-model construction. They should be treated
as verification artefacts, not as substitutes for the formal
argument in the main text. Both the Python script
\texttt{core\_formal\_model\_check.py} and the documented
run-output are shipped with the paper; the script's results are
stable across the Python versions tested (3.10, 3.11, 3.12) and
are independent of any third-party library.

\subsubsection{Table A: Sixteen interpretations of the Boolean core}
\label{sec:verification-table}

The Boolean fragment consists of the four predicates
$\operatorname{Causal}$, $\operatorname{Custom}$, $\operatorname{Bel}$,
$\operatorname{Just}$ evaluated on a single $(b,c)$ pair. There are
$2^4=16$ interpretations. The columns are: truth values of the
four predicates; truth of $M_0$ (i.e.\ M$_{0a}\land$M$_{0b}$); truth
of $T_1$; truth of $T_2$. ``\checkmark'' denotes true; ``$\times$''
denotes false.

\begin{center}
\small
\begin{tabular}{cccc|ccc|c}
\hline
$\operatorname{Causal}$ & $\operatorname{Custom}$ &
$\operatorname{Bel}$ & $\operatorname{Just}$ & $M_0$ & $T_1$ & $T_2$ & note \\
\hline
0 & 0 & 0 & 0 & \checkmark & \checkmark & \checkmark & \\
0 & 0 & 0 & 1 & \checkmark & \checkmark & \checkmark & \\
0 & 0 & 1 & 0 & \checkmark & \checkmark & \checkmark & \\
0 & 0 & 1 & 1 & \checkmark & \checkmark & \checkmark & \\
0 & 1 & 0 & 0 & \checkmark & \checkmark & \checkmark & \\
0 & 1 & 0 & 1 & \checkmark & \checkmark & \checkmark & \\
0 & 1 & 1 & 0 & \checkmark & \checkmark & \checkmark & \\
0 & 1 & 1 & 1 & \checkmark & \checkmark & $\times$ & \textbf{unique countermodel} \\
1 & 0 & 0 & 0 & $\times$ & \checkmark & \checkmark & M$_0$a \& M$_0$b violation \\
1 & 1 & 0 & 0 & $\times$ & \checkmark & \checkmark & M$_0$b violation \\
1 & 0 & 1 & 0 & $\times$ & \checkmark & \checkmark & M$_0$a violation \\
1 & 0 & 0 & 1 & $\times$ & $\times$ & \checkmark & \\
1 & 0 & 1 & 1 & $\times$ & $\times$ & \checkmark & \\
1 & 1 & 0 & 1 & $\times$ & $\times$ & \checkmark & \\
1 & 1 & 1 & 0 & \checkmark & \checkmark & \checkmark & \\
1 & 1 & 1 & 1 & \checkmark & $\times$ & $\times$ & \\
\hline
\end{tabular}
\end{center}

Two facts are worth highlighting.
\begin{enumerate}
\item The countermodel to $T_1\land M_0\models T_2$ is
\emph{unique}: of the 16 interpretations, exactly one
($(0,1,1,1)$) falsifies the entailment. The finite Boolean
fragment is therefore as informative as the general
Lemma~\ref{lem:reduct-invariance} for the single-pair case.
\item The bridge collapse of
Proposition~\ref{prop:bridge-collapse} is sharp: adding $B_0$
eliminates \emph{all} falsifying interpretations of
$T_1\land B_0\land M_0\models T_2$ (verified by adding the
$B_0$ column to the same 16-row table).
\end{enumerate}
The computational evidence is therefore a finite-validation
artefact, not a substitute for the general proof, and is consistent
with all three propositions of \S\ref{sec:prop2} as well as with the
characterization $(\star)$ of
\S\ref{sec:bridge-collapse}.

\subsection{Scope of the Formal Result}
\label{sec:formal-scope}

The formal result is relative to the language, ontology, and
interpretive choices specified above. It does not establish that:
\begin{enumerate}
\item no richer first-order language could represent the relevant
      distinction;
\item modal logic is philosophically superior to grounding or
      justification logic;
\item Stoic katalepsis is identical to any contemporary theory of
      reliabilism;
\item Hume denies every form of epistemic justification, or that
      the ``strong exclusion reconstruction'' (H-I) is the uniquely
      correct reading of Hume's appeal to custom;
\item the formal reconstruction uniquely determines the historical
      interpretation of the Stoic or Humean texts.
\end{enumerate}
The claim is narrower: an extensional reconstruction at the
specified level of granularity, on the specified class of
admissible enrichments $\mathcal K$, can preserve material content
while leaving the explanatory and modal relations semantically
underdetermined.

\begin{quote}
This article does not show that first-order logic cannot represent
grounding, modality, or epistemic normativity. It shows that, for
the explicitly specified extensional signature $L_0$ and class of
admissible reconstructions $\mathcal K$, the intended explanatory
relation between custom-production and the absence of rational
warrant is not fixed by the $L_0$-reduct; representing it requires
additional structure whose philosophical interpretation remains
independently contestable.
\end{quote}


\subsection{Stoic Application: Encoding Sensitivity}
\label{sec:encoding-sensitivity}

The Stoic material admits two levels of formalization, and the
robustness of the negative result should be tested against both. Let
$L_0^A$ be the minimal encoding of \S\ref{sec:stoic-skeleton}, which
registers only the kataleptic predicate $Kat(i)$. Let $L_0^B$ be a
provenance-sensitive encoding that decomposes the kataleptic condition
into a veridicality clause, a source-match clause, and a modal
source-reliability clause:
\begin{equation}
Kat(i) \;\longleftrightarrow\; Veridical(i) \wedge SourceMatch(i)
\wedge \Box_{\mathrm{src}}\,\neg FalseSource(i).
\label{eq:kat-B}
\end{equation}
Formula~\eqref{eq:kat-B} is a formal decomposition adopted for the
present reconstruction; it is not presented as a historical Stoic
axiom. The two encodings differ in where the modal content is placed:
$L_0^A$ leaves the modal clause as an external residue (cf.\
\S\ref{sec:stoic-modal-note}), whereas $L_0^B$ internalizes it by
adding modal vocabulary to the signature.

The structural fact that matters is the same in both encodings: the
modal clause is not a logical consequence of the minimal extensional
predicates alone. In $L_0^A$ it is not even expressible and is recorded
as a residue; in $L_0^B$ it is expressible only because
$\Box_{\mathrm{src}}$ is a new semantic resource beyond the
extensional base. Under either encoding, the minimal extensional
reconstruction does not by itself determine the modal content of the
Stoic definition.

Whether the two encodings yield the \emph{same} underdetermination
result is a question about the reconstruction, and it was tested
computationally (script \texttt{encoding\_sensitivity\_check.py} in
the reproducibility package; frozen output
\texttt{encoding\_sensitivity\_output.txt}). The test enumerates,
over the finite Boolean fragment, every base assignment
$E=(Kat,Veridical,SourceMatch,FalseSource(i))$ and every modal
resource $(R(i,i),R(i,j),FalseSource(j))$, and asks, for each
encoding, whether the modal content $M$ is fixed by $E$ alone. The
outcome is encoding-sensitive in degree, robust in existence:

\begin{itemize}
\item Under $L_0^A$ the modal content $M$ is undetermined for all
$16$ base assignments: for every $E$, both $M=0$ and $M=1$ are
realized by suitable modal resources.
\item Under $L_0^B$ the modal content is undetermined for $6$ of the
$10$ admissible base assignments and determined for the remaining
$4$ --- the kataleptic assignment ($Kat=1$, $Veridical=1$,
$SourceMatch=1$) and its direct negation ($Kat=0$, $Veridical=1$,
$SourceMatch=1$) --- where the determination is produced by the
decomposition axiom itself, i.e.\ by the modal resource written into
the theory.
\end{itemize}

The two encodings therefore do not yield the same underdetermination
result: the strength of the residue is encoding-dependent. This is the
finding, not a defect: it exhibits the hermeneutic dependence of the
formalization. The modal clause is a residue under every encoding; the
only question is whether the residue is left external ($L_0^A$) or
internalized at the cost of adding modal vocabulary to the theory
($L_0^B$). Under both encodings the general claim of this paper is
confirmed: the extensional base alone does not fix the modal content,
and representing it requires a semantic resource beyond the
extensional base.

\subsection{Expressive Enrichments: The Minimal-Enlargement Benchmark}
\label{sec:e0e1e2}

\S\ref{sec:richer-languages} records the three enrichments --- modal,
justification-logic, and grounding --- that can represent what $L_0$
leaves underdetermined. The philosophical content of that section is
sharpened by a minimal-enlargement benchmark, which separates what each
enrichment \emph{adds to the language} from what it \emph{establishes
about the represented relation}.

\begin{center}
\begin{tabular}{@{}lccc@{}}
\toprule
 & \emph{Modality} & \emph{Explanatory direction} & \emph{Hyperintensional distinction}\\
\midrule
E0: extensional $L_0$ & --- & --- & ---\\
E1: $L_0 + G$ (primitive $G$) & --- & $\checkmark^{*}$ & limited\\
E2: $L_0 + G + \Gamma$ & optional & $\checkmark$ & $\checkmark$\\
\bottomrule
\end{tabular}
\end{center}

\noindent $^{*}$A primitive relation $G(x,y)$ formally represents a
binary relation; that it is \emph{grounding} --- and not merely a
relation bearing that name --- is a further semantic claim that
requires the theoretical constraints $\Gamma$ of E2.

The benchmark makes explicit the paper's central distinction:
\textbf{representability is not adequate representation}. E1 makes the
explanatory relation \emph{nameable}; E2 makes it \emph{characterized}.
The negative result of \S\ref{sec:prop2} concerns the step from E0 to
E1 --- the enrichment is not recoverable from $L_0$ alone --- and the
step from E1 to E2 shows that naming is not yet theorizing: the
constraints $\Gamma$ (irreflexivity, asymmetry, transitivity), which
are contested in the grounding literature, are themselves interpretive
choices rather than consequences of the signature.

\subsection{Gate 1.5: Non-Triviality / Recoverability Check}
\label{sec:gate15}

The central theorem is exposed to the objection that it merely records
that a symbol absent from $L_0$ is absent from $L_0$. The
non-triviality of the result is fixed by the following ten-point check,
which the formal section must pass in full before any historical claim
is allowed to rest on it.

\begin{enumerate}
\item[T1.] $T_G$ is explicitly defined: $L_G$, $T_G$, $\mathrm{Models}(T_G)$.
\item[T2.] At least two enriched structures $\mathcal M_1^G, \mathcal M_2^G \models T_G$ exist.
\item[T3.] Same $L_0$-reduct: $\mathcal M_1^G\!\upharpoonright L_0 = \mathcal M_2^G\!\upharpoonright L_0$.
\item[T4.] Different $G$: $G^{\mathcal M_1} \neq G^{\mathcal M_2}$.
\item[T5.] Both structures satisfy the grounding-theoretic constraints $\Gamma$.
\item[T6.] The result is statable independently of the fact that $G \notin L_0$ (non-trivial).
\item[T7.] The definability application is formally valid: implicit-definability definition; absence of explicit definition.
\item[T8.] The E1/E2 comparison is complete: representability vs.\ adequacy.
\item[T9.] The sensitivity test across at least two reasonable encodings ($L_0^A$, $L_0^B$) is performed.
\item[T10.] Machine verification is separated from proof: computational checks are reproducibility artefacts, not substitutes for the semantic argument.
\end{enumerate}

The ten items are verified individually; the complete check is
recorded in Table~\ref{tab:gate15}.

\begin{table}[ht]
\centering
\begin{tabular}{@{}lp{3.6cm}p{5.6cm}c@{}}
\toprule
\textbf{Item} & \textbf{Requirement} & \textbf{Discharged by} & \textbf{Status}\\
\midrule
T1 & $T_G$ explicitly defined: $L_G$, $T_G$, $\mathrm{Models}(T_G)$
& \S\ref{sec:enrichment}--\S\ref{sec:admissible}: $L^+$, axioms
(O$_1$)--(O$_3$), admissible class $\mathcal K$
& $\checkmark$\\
T2 & Two enriched structures $M_1^G, M_2^G \models T_G$ exist
& Proposition~\ref{prop:underdet} construction; script
\texttt{gate15\_check.py}
& $\checkmark$\\
T3 & Same $L_0$-reduct: $M_1^G\!\upharpoonright L_0 = M_2^G\!\upharpoonright L_0$
& By construction in Proposition~\ref{prop:underdet};
Lemma~\ref{lem:reduct-invariance}; script \texttt{gate15\_check.py}
& $\checkmark$\\
T4 & Different $G$: $G^{M_1} \neq G^{M_2}$
& By construction in Proposition~\ref{prop:underdet}; script
\texttt{gate15\_check.py}
& $\checkmark$\\
T5 & Both structures satisfy the constraints $\Gamma$
& Remark below (the constructed pair satisfies
$\Gamma$ in both members); script \texttt{gate15\_check.py}
& $\checkmark$\\
T6 & Result statable independently of $G \notin L_0$ (non-trivial)
& Proposition~\ref{prop:stability} (stability under arbitrary
$L_0$-theories)
& $\checkmark$\\
T7 & Definability application formally valid
& Definitions~2.7--2.8, Proposition~\ref{prop:definability},
\S\ref{sec:definability-qual}
& $\checkmark$\\
T8 & E1/E2 comparison complete: representability vs.\ adequacy
& \S\ref{sec:e0e1e2} (minimal-enlargement benchmark)
& $\checkmark$\\
T9 & Sensitivity test across $L_0^A$, $L_0^B$ performed
& \S\ref{sec:encoding-sensitivity}; script
\texttt{encoding\_sensitivity\_check.py}
& $\checkmark$\\
T10 & Machine verification separated from proof
& \S\ref{sec:verification}: the scripts are reproducibility
artefacts; the general argument is
Lemma~\ref{lem:reduct-invariance}
& $\checkmark$\\
\bottomrule
\end{tabular}
\caption{Gate 1.5 --- non-triviality / recoverability check, 10/10.}
\label{tab:gate15}
\end{table}

Items T1, T6, T7, T8 and T10 are discharged by proof or definition
in the manuscript at the locations cited in the table. Items T2--T5
are verified computationally by \texttt{gate15\_check.py} (frozen
output \texttt{gate15\_output.txt} in the reproducibility package),
which reconstructs the Proposition~\ref{prop:underdet} pair and
asserts the admissibility axioms (O$_1$)--(O$_3$), the equality of the
$L_0$-reducts, the difference of the $G$ extensions, and the
structural constraints $\Gamma$ on both members. Item T9 is
discharged by \texttt{encoding\_sensitivity\_check.py}: the outcome is
encoding-sensitive in degree ($L_0^A$: undetermined for all 16 base
assignments; $L_0^B$: undetermined for 6 of 10 admissible assignments)
and robust in existence --- under both encodings the extensional base
alone fails to fix the modal content. The general thesis survives; the
encoding-dependence of the strength of the residue is recorded as a
finding (hermeneutic dependence), not concealed.

Two remarks complete the check. First, the model pair of
Proposition~\ref{prop:underdet} --- $G_1$ containing
$\langle \mathrm{custFact}(b), \mathrm{nonjustFact}(b,c)\rangle$ and
$G_2$ empty --- satisfies the standard structural constraints $\Gamma$
(irreflexivity, asymmetry, transitivity) in both members: neither
relation contains a reflexive pair, neither contains a pair and its
converse, and neither contains a non-trivial chain. This is verified
by \texttt{gate15\_check.py} (item T5). The underdetermination
therefore does not trade on a pathological choice of $G$, even though
the admissible class $\mathcal K$ deliberately imposes no structural
principles on $G$ (condition 4 of \S\ref{sec:admissible}). Second, T10
is a standing discipline of \S\ref{sec:verification}: the finite
Boolean enumeration, the model-pair script, and the
encoding-sensitivity enumeration are reproducibility artefacts, and
the general arguments are supplied by Lemma~\ref{lem:reduct-invariance}
and the propositions cited in Table~\ref{tab:gate15}. The section does
not proceed to historical inference until all ten items are
discharged.

\subsection{Hyperintensionality: The Four-Layer Claim}
\label{sec:hi-layers}

The formal section can be read as a single four-layer claim about
hyperintensionality. The layers are labelled HI1--HI4 to avoid
confusion with the Humean readings H1--H3 of \S\ref{sec:hume-psr}.

\begin{enumerate}
\item[HI1.] \emph{Literature.} Grounding and explanation are widely
treated in the contemporary literature as hyperintensional phenomena
(SEP ``Metaphysical Grounding''; Fine 2012; Schnieder 2011).
\item[HI2.] \emph{Formal specification.} The extensional language
$L_0$ does not contain explanatory dependence as a primitive semantic
resource (\S\ref{sec:l0}).
\item[HI3.] \emph{Model-theoretic result.} Holding the $L_0$-reduct
fixed, the explanatory enrichment can vary: two admissible structures
with identical $L_0$-reducts realize different grounding readings
(Proposition~\ref{prop:underdet}).
\item[HI4.] \emph{Philosophical interpretation.} The extensional
material profile does not by itself determine the explanatory
structure of the target interpretation.
\end{enumerate}

The chain runs Literature $\to$ Formalization $\to$ Theorem $\to$
Interpretation. The point of separating the layers is to block the
leap from ``grounding is hyperintensional'' to ``first-order logic
cannot express grounding'': HI1 is a claim about the literature, HI2
about a specified signature, HI3 a theorem about that signature, and
HI4 an interpretation of the theorem. The present result occupies
layers HI2--HI4; it neither needs nor asserts the unrestricted version
of HI1 as a theorem about the expressive power of first-order logic.
This is the formal counterpart of the distinction drawn in
\S\ref{sec:e0e1e2} between representability and adequate
representation, and of the modal/grounding distinction
$\Box(Custom(b) \to \neg Just(b,c)) \neq
G(\mathrm{custFact}(b), \mathrm{nonjustFact}(b,c))$ of
\S\ref{sec:richer-languages}.

% =====================================================================
\section{Hume and the Principle of Sufficient Reason: A Qualification
Required}
\label{sec:hume-psr}

A widespread reading says ``Hume rejects the metaphysical PSR'';
another says ``Hume does not reject the PSR; he merely re-describes the
question''. Neither is correct on its own. The verifiable facts are
these.

\begin{enumerate}
\item[(a)] Hume denies that the causal maxim is rationally demonstrable.
In Treatise 1.3.3, the maxim ``whatever begins to exist must have a
cause of existence'' is said to be ``neither intuitively nor
demonstrably certain'': its contrary is conceivable, and our assent to
it arises from observation and experience, i.e., from custom (T 1.3.3;
SBN 78--82). The denial is therefore epistemic: the maxim is not
rationally grounded.
\item[(b)] Hume explicitly denies absolute or metaphysical necessity in
causal connection. T 1.3.14.35: ``there is no absolute nor metaphysical
necessity, that every beginning of existence shou'd be attended with
such an object'' (SBN 172).
\item[(c)] Hume's positive account is psychological. T 1.3.14.20:
necessity ``is nothing but an internal impression of the mind, or a
determination to carry our thoughts from one object to another'' (SBN
164--65). The same thesis is stated in the Enquiry: ``When we look
about us towards external objects, and consider the operation of
causes, we are never able, in a single instance, to discover any power
or necessary connexion'' (E 7.6; SBN 63). And the famous remark that
``reason is nothing but a wonderful and unintelligible instinct in our
souls'' (T 1.3.16.9 SBN 178--9) makes the process-based character of
the account explicit: belief formation is a mechanism, not an
inference.
\end{enumerate}

Hence: Hume does not ``refute'' the PSR in its causal form, and he does
not leave it untouched; he removes it from the level of demonstration
and relocates it at the level of custom. The reading ``Hume denies the
PSR'' is false if it means he argued against the principle's content;
the reading ``Hume never questions it'' is false if it means he
accepted it as a rational principle. The correct formulation is: the
causal maxim is, for Hume, a natural psychological regularity
masquerading as a rational necessity.

\subsection{The four-level Humean analysis}
\label{sec:hume-four-levels}

The formal treatment of \S\ref{sec:formalization} operates on several
levels of claim at once. To keep the formal result from silently
importing a contested interpretive thesis, we record explicitly the
four-level decomposition on which it relies, tagging each level with
its textual anchoring and its standing in the literature (see also
\S\ref{sec:hume-historical} of the formal section).

\begin{center}
\begin{tabular}{p{0.18\linewidth} p{0.42\linewidth} p{0.32\linewidth}}
\hline
\textbf{Level} & \textbf{Claim} & \textbf{Status} \\
\hline
Psychological-genetic &
Causal expectations are formed by custom/habit, not by demonstrative
reason. &
Directly defensible from Hume's text (T 1.3.14.20, SBN 164--65;
E 7.6, SBN 63; T 1.3.16.9, SBN 178--9). \\
Epistemic &
These beliefs are not demonstratively justified. &
Defensible, but the term ``justification'' must be unpacked
(\S\ref{sec:hume-readings}). \\
Normative exclusion &
Custom-production excludes rational justification. &
Contested contemporary reconstruction; alternative readings exist
(H-II, naturalistic; Della Rocca 2010 on PSR as demotion vs.\
rejection; Pruss 2006 on Hume as evaporationist). \\
Grounding &
The custom-production relation \emph{grounds} the absence of
justification. &
Contemporary regimentation in the sense of Fine (2012) and
Schnieder (2011); not directly attributable to Hume. \\
\hline
\end{tabular}
\end{center}

We do not attribute a contemporary grounding relation, nor a fully
articulated theory of epistemic justification, to Hume. We use
grounding only to regiment one strong contemporary reconstruction of
the explanatory force that an interpreter may take Hume's appeal to
custom to possess; alternative reconstructions remain live.

Two further points. First, the term ``Principle of Sufficient Reason''
(\emph{principium rationis sufficientis}) is Leibnizian
(\emph{Monadologie} \S32; 1714); Hume never uses it, and
retro-application requires care. Second, the literature on Hume and the
causal maxim is substantial: on the maxim itself, see Garrett 1997, 24,
66--69; on the projectivist reading of causal necessity, Beebee 2006;
on the Enquiry context, Millican 2002, 27--65. The contemporary revival
of PSR debates (Della Rocca 2010; Pruss 2006) has made the question of
what exactly Hume denied once more a live interpretive issue; our claim
(a)--(c) is deliberately minimal and textually anchored.

% =====================================================================
\section{The Transmission of Stoic Epistemology to Early Modern Europe:
A Verifiable Route}
\label{sec:transmission}

\subsection{Sextus Empiricus}

The kataleptic impression and the criterion debate are documented most
fully in Sextus Empiricus, \emph{Adversus Mathematicos} VII: the
criterion discussion at M 7.24--29, the definitional texts at M 7.248
and 7.402, and the ``clear, evident, striking''
(\emph{enarges}, \emph{ektypos}, \emph{plektike}) marks at M 7.257--58
(cf.\ Bury 1935, 7.402--8). The printing history is now well documented
(Floridi 2002; Popkin 1979, 18): the first Latin translation of the
\emph{Hypotyposes} was published in 1562 by Henri Estienne (Geneva);
the Latin \emph{Adversus Mathematicos}, together with a second edition
of the Estienne translation, was published in 1569 by Gentian Hervet
(Paris: Martin Juvenis; Antwerp: Plantin). The Greek \emph{editio
princeps} of Sextus' \emph{Opera} appeared only in 1621 (Geneva:
Chou\"et). The common claim that ``Hervet translated Sextus in 1562 and
1569'' therefore conflates Estienne with Hervet and should be
corrected.

\subsection{Diogenes Laertius}

The main doxographical source for Zeno's and Chrysippus' epistemology
(the definitions of \emph{phantasia}, \emph{phantasia kataleptike},
\emph{katalepsis}, \emph{epist\=em\=e} at \emph{Vitae Philosophorum}
7.46--54; the criterion at 7.54 [Dorandi 2013]; cf.\ Bobzien 2003,
85--123) reached early modern Europe through Ambrogio Traversari's
Latin translation, completed in manuscript by 1433; the Latin
\emph{editio princeps} was printed in Rome 1472 (Georgius Lauer); the
celebrated Venice 1475 (Nicolaus Jenson) is the second printed edition;
the Greek \emph{editio princeps} appeared in Basel 1533 (Froben)
(Dorandi 2013, 15; ISTC id00219000/id00220000).

\subsection{Cicero and the Neostoic Reception}

Cicero's \emph{Academica} is the Latin source for Zeno's fist analogy
--- open hand: impression; fingers bent: assent; fist: katalepsis; the
other hand grasping the fist: \emph{epist\=em\=e}, available to the
sage alone (\emph{Lucullus} 145 [Brittain 2006]; SVF 1.66) --- and for
the criterion debate in Latin. The \emph{Academica} entered print
early: \emph{editio princeps} Rome 1471 (Sweynheym \& Pannartz), with
the Aldine edition of 1523 establishing the standard text (Hunt 1998).
But the decisive early modern channel for Stoic doctrine as a system
was Neostoicism: Justus Lipsius published \emph{De Constantia} (Antwerp:
Plantin, 1584), and in 1604, both the \emph{Manuductio ad Stoicam
philosophiam} and the \emph{Physiologia Stoicorum} (Antwerp:
Plantin-Moretus). The reception of Stoic ethics through Lipsius and du
Vair is well documented (Lagr\'ee 1994; du Vair 1594); what matters for
this note is that the epistemological transmission --- the katalepsis
doctrine --- had to proceed through Sextus, Diogenes Laertius and
Cicero, since the Neostoics' interest was primarily ethical.

\subsection{Hume's Own Access}

For Hume's own acquaintance with these materials, the documented facts
are these. The catalogue of Hume's library (Norton \& Norton 1996)
lists Cicero's \emph{Opera} (Paris, 1740--42, 9 vols.); it does not
list Sextus Empiricus. The 1840 sale catalogue, on which the
reconstruction rests, shows no Sextus (Fosl 1998). However, the 1840
catalogue reflects books sold at auction; it is not identical to the
full library at Hume's death, as books were often given away or sold
earlier. The absence of Sextus is therefore negative evidence, not a
proof of Hume's ignorance. But building a ``Hume owned a Sextus''
narrative on this gap is historiographically irresponsible. The
standard scholarly account remains: Hume's access to Pyrrhonism ran
through Bayle, Montaigne and others (Popkin 1952, 65--81; 1951). The
historiography of ``Hume the Pyrrhonian'' is best grounded in:
(i)~Hume's direct reading of Cicero; (ii)~his mediated reading of
Sextus through the French sceptical tradition.

\subsection{The Medieval Question and the Cartesian Comparison}

Linear chains of the form ``Stoicism $\to$ Aristotelian commentators
$\to$ Scholasticism $\to$ Descartes $\to$ Locke'' are not documented.
The evidence for discontinuity is strong: only seven Latin Sextus
manuscripts circulated in the medieval Latin West (Floridi 2002); the
\emph{Academica} was among the rarer Cicero texts in the Middle Ages,
with Henry of Ghent the first scholastic explicitly to engage it
(Schmitt 1972; 1983; SEP ``Medieval Skepticism''); the term
``scepticus'' entered Latin usage only in the 1430s. Two qualifications
are required, in the spirit of \S\ref{sec:limits}. First, ``no
continuous direct transmission of the doctrine of katalepsis'' does not
mean ``no knowledge of Stoic epistemology at all'': Cicero's
\emph{comprehensio} and Augustine's \emph{Contra Academicos} provided
an indirect, contested channel (Nawar 2022, 185--207). Second, the
relation between Descartes's ``clear and distinct ideas'' and the
kataleptic impression should be studied as discrete responses to a
common skeptical problem --- the criterion problem --- rather than as a
transmission chain; this is the position of the relevant secondary
literature (Frede 1983; Nawar 2022).

Similarly, Locke's critique of innate principles in Essay I.3 (``No
Innate Practical Principles'') is directed, in its explicit form, at
Herbert of Cherbury's \emph{De Veritate} (1624): Herbert is named at
I.3.15 and discussed through I.3.27 (Nidditch ed. 1975, 77--84); Descartes is not
the primary target of that chapter. And the very notion of a unified
``British Empiricism'' is a construction of nineteenth-century
historiography (Norton 1981). The moral of this section: the history
of philosophy should be written from the history of printed texts, not
from speculative chains.

\subsection{Historical-Evidential Method: The Ev0--Ev4 Ladder}
\label{sec:ev-ladder}

The historical discipline of \S\ref{sec:transmission} can be made
explicit as an evidence ladder. The labels are Ev0--Ev4 to avoid
confusion with the expressive-enrichment levels E0/E1/E2 of
\S\ref{sec:e0e1e2}.

\begin{itemize}
\item[Ev0] \emph{Bibliographic presence}: the text exists and is
attested.
\item[Ev1] \emph{Documented access}: an agent or library can be shown
to have had the text available.
\item[Ev2] \emph{Documented citation}: the text is cited or quoted by
the agent.
\item[Ev3] \emph{Textual dependence}: the agent's claims require the
text's specific content.
\item[Ev4] \emph{Influence}: the text shaped the agent's doctrine.
\end{itemize}

Rule: Ev$_n \not\Rightarrow$ Ev$_{n+1}$ in general. Two prohibitions
follow, which the present note observes throughout. First,
\emph{availability is not influence}: that Sextus was printed in 1562
does not show that Hume's doctrine was shaped by him. Second,
\emph{citation is not dependence}: that an author cites a text does
not show that the cited content is load-bearing for the author's
argument. The transmission route of \S\ref{sec:transmission}.1--3 is
established at the level of Ev1--Ev2 (documented access and
citation); Hume's own acquaintance with Sextus is at most Ev1--Ev2 via
Bayle and the French sceptical tradition
(\S\ref{sec:transmission}.4); and the medieval chain is rejected
because it does not reach Ev1 (\S\ref{sec:transmission}.5). No claim
of Ev4 is made anywhere in this note.

% =====================================================================
\section{A Comparative Note: Formal Limits Across Philosophical
Traditions}
\label{sec:comparative}

The theme of this note --- that a formal scheme has limits that must be
stated from within and beyond --- is not peculiar to the Western
analytic tradition. We record three documented parallels and note two
further items for completeness. Each is invoked as an analogy only: no
claim of historical transmission between these traditions and the
Stoic/Humean material is made or implied (cf.\ \S\ref{sec:limits}).

\begin{enumerate}
\item[(a)] \textbf{The catu\c{s}ko\d{t}i.} The Buddhist four-cornered
negation schema --- $P$, $\neg P$, both, neither --- was theorized by
N\=ag\=arjuna as an exhaustive decomposition of assertoric space, and
it has been shown that classical two-valued semantics cannot express
the fourth corner as a distinct position; the schema is naturally
formalized in paraconsistent logics (Priest 2010, 24--54; 2018). One
asymmetry must be stated honestly: the catu\c{s}ko\d{t}i is a case of
expanding a formal scheme --- paraconsistency extends the logic ---
whereas our \S\ref{sec:formalization} result is a case of a scheme's
representational limit. The parallel is therefore not in the kind of
limit but in the practice: the catu\c{s}ko\d{t}i shows that theorizing
about what a scheme can and cannot say --- including the scheme's own
boundary --- is a distinguished non-Western philosophical activity.
Whether N\=ag\=arjuna himself intended the catu\c{s}ko\d{t}i as a
meta-logical reflection on the limits of assertoric space, or whether
this is a modern reconstruction (Priest 2018, ch.~1; Tillemans 1999),
is itself a contested interpretive question. The claim is not that
Stoic or Humean materials require paraconsistency; whatever
N\=ag\=arjuna's intentions, the catu\c{s}ko\d{t}i schema itself --- as
an exhaustive decomposition of assertoric space --- encodes a
boundary-theorizing practice, and so the question ``what cannot be
represented here?'' is not a modern Western invention.

\item[(b)] \textbf{The Xunzian rectification of names (zhengming).}
Xunzi's chapter ``On the Rectification of Names'' (Xunzi 22) addresses
the criterion problem for language: names are correct if fixed by
shared sensory experience (the \emph{tianguan}, ``heavenly faculties'',
common to all humans --- Xunzi 22.5) and by convention (``names have no
intrinsic appropriateness; they are fixed by agreement'' --- Xunzi
22.9; Hansen 1992, ch.~9; Hansen 1983). The structural parallel with
the Stoic criterion debate is striking and has been noted in the
comparative literature (Ro\v{s}ker, SEP ``Epistemology in Chinese
Philosophy''; Kjellberg 1996): both theories face the question of what
fixes the standard by which representations are judged correct. The
differences are equally instructive. At the level of the kind of
standard, the Stoic criterion is source-based (the impression's
provenance) while the Xunzian criterion is intersubjective-conventional.
But the deeper difference is in the functional position each doctrine
occupies within its system. The Stoic katalepsis is an epistemological
doctrine about individual cognitive accuracy: it serves to secure truth
for a single cognitive agent, and its failure mode is epistemic error.
The Xunzian zhengming, by contrast, is an ethical-political doctrine
about the ordering of social life: its purpose is to ensure that
correct naming yields correct action and social coordination (cf.\
Xunzi 22.9; the canonical formula ``if names are not correct, language
will not be in accordance with truth'' is \emph{Analects} 13.3 --- the
Xunzian claim is a development in the same spirit rather than a verbatim
citation), and its failure mode is social disorder. The structural
parallel at the level of ``what fixes the standard?'' thus masks a
difference in the domain and the end to which the standard is put: the
analogy should not be read as implying that the two doctrines occupy
the same functional position in their respective systems. A first-order
encoding of either faces the same underdetermination problem we isolate
in \S\ref{sec:prop2}: the criterion's normative force --- in virtue
of what the standard is the standard --- is not a truth-condition.

\item[(c)] \textbf{Kaozheng and documentary discipline.} The Qing-dynasty
kaozheng (\emph{kao\-ju}; ``evidential research'') movement ---
exemplified by Dai Zhen (\emph{Mengzi ziyi shuzheng}, Preface; in
\emph{Dai Zhen wenji}, Zhonghua Shuju, 1980; cf.\ Elman 1984,
Introduction and ch.~5; 2nd ed.\ 2001) --- made textual verification,
philological evidence and the rejection of unfounded commentary the
methodological core of scholarship. Our \S\ref{sec:transmission} and
\S\ref{sec:limits} defend precisely this discipline for the history of
philosophy: the claim that documentary discipline is a method, not an
ornament, has a systematic analogue in a non-Western tradition that is
still under-cited in analytic historiography. The parallel is in the
methodological attitude --- evidence over speculation --- not in the
object of inquiry, which differs. Again: analogy, not influence.
\end{enumerate}

We also note, for completeness, the Mohist tradition's treatment of the
name--reality (\emph{ming}/\emph{shi}) distinction and the argumentation
chapter ``Lesser Selection'' (\emph{Xiaoqu}) of the \emph{Mozi} (Graham
1978; 1989), and the White Horse paradox of Gongsun Long, which turns
on the gap between a predicate's linguistic ontology and its extension
--- a theme in the vicinity of, though not identical with, the
model-theoretic underdetermination discussed in \S\ref{sec:prop2}
(Hansen 1983, ch.~5). Finally, Wang Yangming's doctrine of
\emph{liangzhi} (``innate knowing''; \emph{Chuanxilu} 2.2; Van Norden,
SEP ``Wang Yangming'') treats the cognitive standard as an inner,
quasi-evidential faculty rather than an external impression or
convention --- structurally closer to the Stoic katalepsis than to the
Xunzian convention, since both posit a source-based standard rather
than an intersubjective one, yet differing in that Wang Yangming's
doctrine is embedded in a moral-spiritual project (the unity of
knowledge and action) rather than primarily in an epistemological
criterion debate. [Provenance note: the chapter and passage numbers
above follow the standard editions cited in the bibliography; where
translations differ (Watson; Knoblock), we rely on the philological
consensus reported in Hansen 1992 and the Stanford Encyclopedia entries
cited. The claim that the Stoic/Xunzian parallel has been noted in the
comparative literature is supported by Kjellberg 1996
(Sextus/Zhuangzi/Xunzi on skepticism) and by the SEP entry
``Epistemology in Chinese Philosophy'' (Ro\v{s}ker); both are cited as
reporting a parallel, not as establishing a historical connection.]

% =====================================================================
\section{Objections and Replies}
\label{sec:objections}

The central result invites a family of objections. They are collected
here following the format \emph{objection -- concession -- distinction
-- response}. No reply introduces new machinery; each restates the
scope fixed in \S\ref{sec:formal-scope} and \S\ref{sec:gate15}.

\medskip
\textbf{Objection 1 (``$G$ is merely a new predicate'').} Adding a
primitive $G$ is something any first-order language can do; the result
therefore shows only that a symbol absent from $L_0$ is absent from
$L_0$.\\
\emph{Concession.} Yes: first-order logic can accommodate a primitive
binary relation $G$.\\
\emph{Distinction.} That $G$ can be \emph{added} does not show that it
is \emph{definable from} $L_0$, nor that the resulting predicate has
the semantic properties characteristic of grounding.\\
\emph{Response.} The claim concerns recoverability from a specified
extensional base, not the impossibility of adding vocabulary; the
non-triviality check of \S\ref{sec:gate15} (item T6) is exactly the
point on which the objection turns.

\medskip
\textbf{Objection 2 (``Modal logic reduces to FOL'').} The standard
translation embeds modal logic in first-order logic; the modal residue
is therefore not a genuine residue.\\
\emph{Concession.} The standard translation is a genuine embedding for
the model-theoretic content of modal logic.\\
\emph{Distinction.} The standard translation preserves \emph{truth in
a model}; it does not preserve explanatory direction. In particular,
$\Box(Custom(b) \to \neg Just(b,c))$ (``it is necessary that \dots'')
is not the grounding claim $G(\mathrm{custFact}(b),
\mathrm{nonjustFact}(b,c))$ (``\dots because \dots'').\\
\emph{Response.} The modal enrichment captures necessity, but not by
itself the explanatory direction; \S\ref{sec:modal-enrichment} states
this explicitly.

\medskip
\textbf{Objection 3 (``Adding $G$ to FOL solves the problem'').} If $G$
can be added, the problem is solved; the paper's own E1 shows this.\\
\emph{Concession.} E1 makes the relation nameable.\\
\emph{Distinction.} Naming is not yet adequacy: a primitive $G$ is not
a grounding theory. The constraints $\Gamma$ of E2 are required for the
relation to be characterized as grounding, and $\Gamma$ is itself an
interpretive choice (\S\ref{sec:e0e1e2}).\\
\emph{Response.} The minimal-enlargement benchmark (E0/E1/E2) is
precisely the experiment that separates ``added vocabulary'' from
``added theory''.

\medskip
\textbf{Objection 4 (``A primitive grounding relation suffices'').}
Grounding may simply be taken as primitive; no theory is needed.\\
\emph{Concession.} A primitive relation can represent a formal
relation.\\
\emph{Distinction.} That the primitive relation \emph{is} grounding ---
and not a merely similar relation --- is a further semantic claim.\\
\emph{Response.} The present result does not deny the possibility of a
primitive $G$; it denies that $G$ is determined by the $L_0$-reduct,
and it records that characterizing $G$ as grounding requires the
semantic constraints of E2.

\medskip
\textbf{Objection 5 (``The result is a semantic triviality'').} The
result states only that a relation absent from the signature is not
determined by that signature.\\
\emph{Distinction.} Trivial lexical absence ($G \notin L_0$) is not
semantic underdetermination ($\mathcal M_1\!\upharpoonright L_0 =
\mathcal M_2\!\upharpoonright L_0$ while $G^{\mathcal M_1} \neq
G^{\mathcal M_2}$). The latter is what
Proposition~\ref{prop:underdet} establishes.\\
\emph{Response.} The Gate 1.5 check (\S\ref{sec:gate15}, item T6) is
designed to make this distinction enforceable rather than asserted.

\medskip
\textbf{Objection 6 (``The Stoic reconstruction is underdetermined'').}
The Stoic material is encoded one way; a different encoding would give
a different result, so the conclusion is an artefact of the choice.\\
\emph{Concession.} The choice of encoding is interpretive and precedes
the formal argument (\S\ref{sec:limits}).\\
\emph{Distinction.} The relevant question is whether the negative
result survives across more than one reasonable encoding.\\
\emph{Response.} The encoding-sensitivity test of
\S\ref{sec:encoding-sensitivity} ($L_0^A$ vs.\ $L_0^B$) is the honest
procedure for answering this; its outcome (robust or sensitive) is
reported rather than assumed (Gate 1.5, item T9).

\medskip
\textbf{Objection 7 (``The Hume reconstruction is controversial'').}
The demotion reading (H2) is not the only interpretation of Hume's
appeal to custom; the result therefore inherits a contested premise.\\
\emph{Concession.} H1 (strong denial) and H3 (evaporation) are live
alternatives in the literature.\\
\emph{Distinction.} The formal result is invariant under the choice
between the exclusion readings modulo the annotations of
\S\ref{sec:hume-readings}; the four-level analysis of
\S\ref{sec:hume-four-levels} is tagged by textual anchoring and
standing.\\
\emph{Response.} H2 is adopted as a methodological interpretation for
reconstruction, not as a settled historical conclusion
(\S\ref{sec:hume-psr}).

\subsection{What the Result Does Not Show: A Negative-Result Matrix}
\label{sec:negative-matrix}

The replies above can be summarized in a matrix that pairs each claim
with what is \emph{not} shown and what \emph{is} shown. The matrix
doubles as an editorial guard: any sentence that asserts a ``not
shown'' claim as a result is a scope violation.

\begin{center}
\begin{tabular}{@{}p{3.6cm}p{4.0cm}p{4.6cm}@{}}
\toprule
\textbf{Claim} & \textbf{What is not shown} & \textbf{What is shown}\\
\midrule
FOL fails & Intrinsic impossibility & $L_0$-relative underdetermination\\
Grounding impossible in FOL & --- & Primitive extension possible\\
$G$ represents grounding & Not by symbol alone & Requires semantic/theoretical constraints $\Gamma$\\
Modal logic solves the explanatory problem & --- & Captures necessity, not explanatory direction\\
Historical reading follows from the formal model & --- & Requires independent documentary evidence (\S\ref{sec:transmission})\\
Richer language cannot represent the distinction & --- & Representation requires an additional semantic resource\\
\bottomrule
\end{tabular}
\end{center}

% =====================================================================
\section{Conclusion}
\label{sec:conclusion}

This note has made three contributions.

\begin{enumerate}
\item \textbf{Technical limit.} A purely extensional first-order
encoding of the justification conditions of Stoic katalepsis and Humean
custom can represent the material content of both positions ---
including the extensionally distinguishable formulations $T_1$ and
$T_2$ (formerly (Exc1) and (Exc2)) --- but cannot fix their explanatory
and modal residue: the ``because'' of the Humean exclusion
(hyperintensional; grounding), the ``could not arise from what is
not'' of the Stoic definition (modal), and the reliability of
belief-forming processes. More precisely, the grounding atom
$G(\mathrm{custFact}(b),\mathrm{nonjustFact}(b,c))$ is \emph{not
implicitly definable} in the extensional signature $L_0$ over the
admissible class of reconstructions $\mathcal K$
(Proposition~\ref{prop:underdet}); the model-pair argument shows the
limit is semantic: the intended reading is not a structural property of
$L_0$-models over $\mathcal K$. A justification operator (Artemov \&
Fitting 2019), a grounding connective (Fine 2012), or a modal operator
is required. This is a singular encoding limit, not a universal failure
claim.
\item \textbf{Interpretive limit.} Hume's relation to the PSR is
neither simple rejection nor silence: he relocates the causal maxim
from the demonstrative level to the level of custom (T 1.3.3; SBN
78--82; T 1.3.14.20; SBN 164--65; T 1.3.14.35; SBN 172). ``Hume never
questions the PSR'' and ``Hume refutes the PSR'' are both false; the
accurate formulation is a demotion, not a refutation. A significant
interpretive counter-thesis --- defended notably by Della Rocca (2010)
and, in a different register, by Pruss (2006) --- argues that this
``demotion'' is a polite fiction: by reducing the PSR to a ``natural
psychological regularity,'' Hume effectively evaporates its normative
status, rendering the distinction between ``demotion'' and
``rejection'' functionally moot. We maintain the demotion formulation
as the more textually anchored reading, since Hume explicitly preserves
the causal maxim as a psychological principle rather than arguing
against its content; but the evaporation reading remains a strong
interpretive alternative, and our claim should be read as a preference
supported by the textual evidence, not as a settled result.
\item \textbf{Historical limit.} The transmission of Stoic epistemology
to early modern Europe runs through documented printed texts: Sextus
(Latin 1562 Estienne; 1569 Hervet; Greek 1621), Diogenes Laertius
(Traversari 1433; Latin \emph{editio princeps} Rome 1472; Venice 1475;
Greek Basel 1533), Cicero's \emph{Academica} (in \emph{Opera omnia},
Rome 1471 [Sweynheym \& Pannartz]; Aldine edition 1523; the fist
analogy at \emph{Lucullus} 145 [Brittain 2006]), and the Neostoic
reception (Lipsius 1604). Speculative medieval chains are not
documented; indirect medieval knowledge via Cicero and Augustine is.
Hume's own access ran through Cicero directly and Sextus through Bayle
(Popkin 1952, 65--81).
\end{enumerate}

Formal tools are useful precisely when their limits are stated; the
history of philosophy requires a documentary discipline that formal
notation cannot supply; and both theses have systematic, independently
developed analogues in the Buddhist, Xunzian and kaozheng traditions
(\S\ref{sec:comparative}).

% =====================================================================
\section{Limits of This Note}
\label{sec:limits}

In the same spirit, the limits of this note are stated explicitly.

\begin{enumerate}
\item We make no claim of in-principle inexpressibility for any
language richer than $L_0$: the negative result is about this
signature, at this granularity, under extensional semantics, over the
admissible class $\mathcal K$.
\item We do not claim that justification logic, grounding, or modal
logic ``solves'' the philosophical problems; we claim only that they
make the relevant distinction representable.
\item The comparative section (\S\ref{sec:comparative}) is analogical;
none of its claims depends on a historical connection between the
traditions compared, and none asserts one.
\item All historical claims rest on the printed-text history summarized
in \S\ref{sec:transmission}; where the secondary literature is the sole
support (e.g., the medieval manuscript census), this is stated, and the
sources are given.
\item ``No evidence of continuous direct transmission''
(\S\ref{sec:transmission}.5) is a negative-evidence claim; it is
recorded as such and is not asserted to be a proof of non-transmission.
\item Quotations follow the editions listed under ``Citations and
Editions''; paragraph numbers for Hume follow the Norton \& Norton
(2000) and Beauchamp (1999) editions, with Selby-Bigge/Nidditch (SBN)
pagination given in every case.
\item The formal encoding in \S\ref{sec:l0} is itself an interpretive
choice. Different readings of the Stoic and Humean sources might yield
different signatures or axioms; our result is relative to the encoding
chosen, and the choice of encoding is a hermeneutic decision that
precedes the formal argument.
\item The richer frameworks of \S\ref{sec:richer-languages} ---
grounding, justification logic, process reliabilism --- are invoked as
contemporary regimentations of the relevant claims; no claim is made
that Hume or the Stoics held these theories, and their vocabulary is
applied in a deliberately anachronistic, analytic sense.
\end{enumerate}

% =====================================================================
\section{Appendix: Definitions}
\label{sec:appendix}

\begin{center}
\begin{longtable}{p{0.20\linewidth} p{0.45\linewidth} p{0.28\linewidth}}
\hline
\textbf{Term} & \textbf{Definition} & \textbf{Source} \\
\hline
katalepsis &
Stoic ``cognition'': the grasping of a kataleptic impression; the
criterion of truth for the Stoics &
DL 7.46, 7.54 [Dorandi 2013]; translation ``cognition'': Long \&
Sedley 1987, 40; SEP ``Stoicism'' \S3.7 \\
\hline
phantasia kataleptike &
Kataleptic impression: ``arising from what is, stamped and impressed
in accordance with what is, of such a kind as could not arise from
what is not'' &
DL 7.46 [Dorandi 2013]; M 7.248 [Bury 1935]; 7.402 [Bury 1935] \\
\hline
Criterion problem &
The question what standard fixes which impressions are true; the
regress of ``criterion of the criterion'' &
Sextus M 7.24--29; PH I.164--169 (169 = diallelos); M 7.440--446 \\
\hline
Assent (sunkatathesis) &
The Stoic act of approving an impression; the middle stage between
impression and cognition &
DL 7.46--54; Cicero, \emph{Acad.} 2.37--39, 2.66--67 \\
\hline
katalepsis / episteme &
The Stoic hierarchy: katalepsis is the assent to a kataleptic
impression (sunkatathesis kataleptikei phantasiai); episteme is
secure cognition (katalepsis asphales) &
M 7.151--152 [Bury 1935]; DL 7.47 [Dorandi 2013] \\
\hline
PSR &
\emph{Principium rationis sufficientis}: Leibnizian term; every fact
has a sufficient reason &
Leibniz, \emph{Monadologie} \S32 (1714) \\
\hline
FOL &
First-order logic: extensional, truth-functional model-theoretic
semantics &
Standard; \S\ref{sec:l0} \\
\hline
$L_0$ &
The minimal many-sorted extensional language of the formal section:
sorts $I$ (impressions), $\mathit{Cont}$ (contents), $B$ (cognitive
episodes); predicates $\operatorname{Kat}$, $\operatorname{Rep}$,
$\operatorname{Grasp}$, $\operatorname{Assent}$,
$\operatorname{Bel}$, $\operatorname{Causal}$, $\operatorname{Custom}$,
$\operatorname{Just}$, $\operatorname{StoicEp}$ &
\S\ref{sec:l0} \\
\hline
$L^+$ &
The enriched language extending $L_0$ with reified fact-objects
$\mathit{Fact}\subseteq\mathit{Cont}$, the functions
$\operatorname{custFact}:B\to\mathit{Fact}$ and
$\operatorname{nonjustFact}:B\times\mathit{Cont}\to\mathit{Fact}$,
the grounding relation $G\subseteq\mathit{Fact}\times\mathit{Fact}$,
and the obtaining monad $\operatorname{Obtains}$ &
\S\ref{sec:enrichment} \\
\hline
$\mathcal K$ &
The class of admissible reconstructions: $L_0$-structures satisfying
$M_0$ and the sortal and reification axioms $O_1$--$O_3$, with no
bridge axiom and no structural principles on $G$& \S\ref{sec:admissible} \\
\hline
Grounding &
Hyperintensional explanatory relation: $\phi < \psi$, ``$\phi$ grounds
$\psi$''; in this paper regimented on reified fact-objects
($G \subseteq \mathit{Fact}\times\mathit{Fact}$)& Fine 2012; Correia \& Schnieder 2012; \S\ref{sec:enrichment} \\
\hline
Implicit definability &
A relation $R$ is implicitly definable in $L$ over a class $\mathcal K$
iff any two admissible structures sharing an $L$-reduct agree on $R$& Beth 1953; \S\ref{sec:definability-qual} \\
\hline
Custom &
Humean habit: determination of the mind to pass from one object to its
usual attendant &
T 1.3.14.20; E 5 \\
\hline
Minimal dependency skeleton &
The Stoic reconstruction (S): $\forall b\,(\operatorname{StoicEp}(b)
\to \exists i\exists c\,[\operatorname{Kat}(i)\land
\operatorname{Rep}(i,c)\land \operatorname{Grasp}(b,i)\land
\operatorname{Assent}(b,c)])$; deliberately excludes truth, guarantee
and the modal clause &
\S\ref{sec:stoic-skeleton} \\
\hline
Four Humean levels &
Psychological-genetic / epistemic / normative exclusion / grounding;
see \S\ref{sec:hume-four-levels} &
\S\ref{sec:hume-four-levels} \\
\hline
catu\c{s}ko\d{t}i &
Buddhist four-cornered negation: $P$, $\neg P$, both, neither &
Priest 2010 \\
\hline
zhengming (\emph{zheng ming}) &
Xunzian rectification of names: fixing correct usage by sensory
agreement and convention &
Xunzi 22.5, 22.9; Hansen 1992 \\
\hline
kaozheng (\emph{kao\-ju}) &
Qing evidential scholarship: philological verification as method &
Elman 1984, Introduction and ch.~5 \\
\hline
\end{longtable}
\end{center}

\section*{Citations and Editions}

\begin{itemize}
\item Hume, \emph{Treatise} (T): Norton \& Norton (eds.), Oxford
University Press, 2000; paragraphs cited as T 1.3.3 etc.;
Selby-Bigge/Nidditch (SBN) pagination in parentheses.
\item Hume, \emph{Enquiry} (E): Beauchamp (ed.), Oxford University
Press, 1999; paragraphs cited as E 7.6 etc.; SBN pagination in
parentheses.
\item Sextus Empiricus: Primary source editions: Latin translations
--- \emph{Hypotyposes}, tr.\ Henri Estienne, 1562;
\emph{Adversus Mathematicos}, tr.\ Gentian Hervet, 1569; Greek
\emph{Opera}, Chou\"et, 1621. Modern translation used: Bury (tr.), Loeb Classical
Library, 1933--49.
\item Diogenes Laertius: Hicks (tr.), Loeb Classical Library, 1925;
Dorandi (ed.), Cambridge University Press, 2013.
\item Cicero, \emph{Academica}: Brittain (tr.), Hackett, 2006;
\emph{Lucullus} paragraph numbers (Plasberg).
\item Dai Zhen: \emph{Mengzi ziyi shuzheng} (\emph{Men\-gzi zi\-yi
shu\-zheng}), in \emph{Dai Zhen wenji} (\emph{Dai Zhen wen\-ji}),
Beijing: Zhonghua Shuju, 1980.
\item von Arnim, H. (ed.). \emph{Stoicorum Veterum Fragmenta} (SVF),
vol.~1. Leipzig: Teubner, 1905.
\end{itemize}

\section*{Provenance Table}

\begin{center}
\begin{longtable}{p{0.14\linewidth} p{0.48\linewidth} p{0.32\linewidth}}
\hline
\textbf{\S} & \textbf{Claim} & \textbf{Source} \\
\hline
Abstract &
katalepsis = ``cognition'' (Long \& Sedley 1987, 40) &
Long \& Sedley 1987, 40; cf.\ SEP ``Stoicism'' \S3.7 [P-01] \\
\hline
\S\ref{sec:stoic-modal-note} &
Kataleptic definition with modal clause &
DL 7.46 [Dorandi 2013]; M 7.248 [Bury 1935]; 7.402 [Bury 1935]
[P-03] \\
\hline
\S\ref{sec:appendix} &
katalepsis = sunkatathesis kataleptikei phantasiai; episteme =
katalepsis asphales &
M 7.151--152 [Bury 1935]; DL 7.47 [Dorandi 2013] [P-03b] \\
\hline
\S\ref{sec:prop1} &
$T_2\land M_0\models T_1$; $T_1\land M_0\not\models T_2$; unique
countermodel $(0,1,1,1)$ &
Formal argument; Table A (16 Boolean rows) [P-04] \\
\hline
\S\ref{sec:prop1} &
Bridge collapse: $T_1\land M_0\land B_0\models T_2$ (0 falsifiers) &
Formal argument; machine-checked [P-04b] \\
\hline
\S\ref{sec:enrichment} &
Grounding hyperintensional; regimented on reified fact-objects &
Fine 2012; Correia \& Schnieder 2012; Schnieder 2011 [P-05] \\
\hline
\S\ref{sec:prop2} &
Grounding atom not implicitly definable in $L_0$ over $\mathcal K$ &
Formal argument; Lemma (reduct invariance); Beth 1953 [P-05b] \\
\hline
\S\ref{sec:definability-qual} &
Explicit vs.\ implicit definability; Beth anchor& Beth 1953 [P-05c] \\
\hline
\S\ref{sec:richer-languages} &
Process reliabilism &
Goldman 1979, 1--23; 1986 [P-05] \\
\hline
\S\ref{sec:richer-languages} &
Justification logic &
Artemov 2008; Artemov \& Fitting 2019 [P-06] \\
\hline
\S\ref{sec:hume-psr} &
T 1.3.3; SBN 78--82 &
Hume, T (Norton \& Norton 2000); SBN [P-07] \\
\hline
\S\ref{sec:hume-psr} &
T 1.3.14.20; SBN 164--65 &
Hume, T [P-07] \\
\hline
\S\ref{sec:hume-psr} &
T 1.3.14.35; SBN 172 &
Hume, T [P-07] \\
\hline
\S\ref{sec:hume-psr} &
E 7.6; SBN 63 &
Hume, E (Beauchamp 1999) [P-07] \\
\hline
\S\ref{sec:hume-psr} &
T 1.3.16.9; SBN 178--9 &
Hume, T [P-07] \\
\hline
\S\ref{sec:hume-psr} &
PSR literature &
Garrett 1997; Beebee 2006; Millican 2002; Della Rocca 2010; Pruss
2006 [P-07] \\
\hline
\S\ref{sec:transmission}.1 &
1562 Estienne; 1569 Hervet; Greek 1621 &
Floridi 2002; Popkin 1979, 18 [P-08] \\
\hline
\S\ref{sec:transmission}.1 &
M 7.24--29; M 7.257--58 &
Sextus, M (Bury, Loeb) [P-08] \\
\hline
\S\ref{sec:transmission}.2 &
Traversari 1433; Rome 1472; Venice 1475; Basel 1533 &
Dorandi 2013, 15; ISTC id00219000/id00220000 [P-09] \\
\hline
\S\ref{sec:transmission}.3 &
Fist analogy at \emph{Lucullus} 145 [Brittain 2006] &
Cicero, \emph{Acad.} (Brittain 2006); SVF 1.66 [P-10] \\
\hline
Abstract &
Fist analogy = Long \& Sedley 41B &
Long \& Sedley 1987, section 41B (= Cicero, \emph{Lucullus} 145)
[P-01] \\
\hline
\S\ref{sec:transmission}.3 &
Lipsius 1584; 1604 &
SEP ``Justus Lipsius''; Lagr\'ee 1994 [P-10] \\
\hline
\S\ref{sec:transmission}.4 &
Hume's library: Cicero yes, Sextus no &
Norton \& Norton 1996; Fosl 1998 [P-11] \\
\hline
\S\ref{sec:transmission}.4 &
Bayle-mediated access &
Popkin 1952, 65--81; 1951 [P-11] \\
\hline
\S\ref{sec:transmission}.5 &
Medieval discontinuity; nuances &
Floridi 2002; Schmitt 1972; SEP ``Medieval Skepticism''; Nawar 2022
[P-12] \\
\hline
\S\ref{sec:transmission}.5 &
Locke--Herbert I.3.15--27 &
Locke, \emph{Essay} (Nidditch 1975), 77--84 [P-12] \\
\hline
\S\ref{sec:transmission}.5 &
British Empiricism myth &
Norton 1981 [P-12] \\
\hline
\S\ref{sec:comparative} &
Catu\c{s}ko\d{t}i formalized &
Priest 2010; 2018 [P-13] \\
\hline
\S\ref{sec:comparative} &
Xunzi 22.5; 22.9 &
Xunzi (ctext; Knoblock tr.); Hansen 1992, ch.~9 [P-13] \\
\hline
\S\ref{sec:comparative} &
Kaozheng &
Elman 1984, Introduction and ch.~5; 2nd ed.\ 2001 [P-13] \\
\hline
\S\ref{sec:comparative} &
Mohist/Gongsun Long &
Graham 1978; 1989; Hansen 1983, 140--51 [P-13] \\
\hline
\S\ref{sec:limits} &
Limits stated &
This note [P-15] \\
\hline
\end{longtable}
\end{center}

\noindent SC-ID mapping (CHK-02). Every mandatory SC-ID corresponds to
at least one section covered by the table: SC-001--004 $\to$
Abstract/\"{O}zet/Keywords; SC-005 $\to$ \S\ref{sec:intro}; SC-006--012
$\to$ \S\ref{sec:l0}--\S\ref{sec:richer-languages}; SC-013--014 $\to$
\S\ref{sec:hume-psr}; SC-015--017 $\to$ \S\ref{sec:transmission};
SC-018 $\to$ \S\ref{sec:comparative}; SC-019 $\to$ \S\ref{sec:conclusion};
SC-020--022 $\to$ Provenance/Citations/References; SC-023 $\to$
Appendix: Definitions; SC-024--030 $\to$ \S\ref{sec:transmission}.4,
\S\ref{sec:transmission}.5, Appendix, \S\ref{sec:limits} (optional
items are covered where adopted).

\section*{Open Science Statement}

The initial draft was generated in 2026 with the assistance of an AI
language model used for formalization and bibliography retrieval. All
models, ISTC numbers, SBN citations, and Greek transcriptions were
manually verified against primary sources and critical editions.
The reproducibility package accompanying this manuscript (the formal
model-check script, its frozen output, and the manifest) records the
verification artefacts separately from the mathematical argument.
The human author is responsible for the final content.

\section*{References}

\begin{itemize}
\item Artemov, S. (2008). ``The Logic of Justification.'' \emph{Review
of Symbolic Logic} 1(4): 477--513.
\item Artemov, S. \& Fitting, M. (2019). \emph{Justification Logic:
Reasoning with Reasons}. Cambridge University Press.
\item Beauchamp, T.L. (ed.) (1999). \emph{An Enquiry Concerning Human
Understanding}, by David Hume. Oxford University Press.
\item Beebee, H. (2006). \emph{Hume on Causation}. Routledge.
\item Beth, E.W. (1953). ``On Padoa's Method in the Theory of
Definition.'' \emph{Indagationes Mathematicae} 15: 330--339
(also \emph{Proc. Kon. Ned. Akad. Wetensch.} A56: 330--339).
\item Bobzien, S. (2003). ``Stoic Logic.'' In B.~Inwood (ed.), \emph{The
Cambridge Companion to the Stoics}, 85--123. Cambridge University
Press.
\item Brittain, C. (tr.) (2006). \emph{Cicero: On Academic Scepticism}.
Hackett.
\item Bury, R.G. (tr.) (1933--49). \emph{Sextus Empiricus}, 4 vols.
Loeb Classical Library. (Citations by volume year: ``Bury 1935'' = vol.~II,
\emph{Against the Logicians}, M 7--8.)
\item Cicero. \emph{Academica}. Brittain 2006; \emph{Lucullus} 145
(fist analogy).
\item Correia, F. \& Schnieder, B. (eds.) (2012). \emph{Metaphysical
Grounding: Understanding the Structure of Reality}. Cambridge
University Press.
\item Della Rocca, M. (2010). ``PSR.'' \emph{Philosophers' Imprint}
10(7).
\item Diogenes Laertius. \emph{Vitae Philosophorum} VII. Hicks 1925;
Dorandi 2013.
\item Dorandi, T. (ed.) (2013). \emph{Diogenes Laertius: Lives of
Eminent Philosophers}. Cambridge University Press.
\item Elman, B.A. (1984). \emph{From Philosophy to Philology:
Intellectual and Social Aspects of Change in Late Imperial China}.
Council on East Asian Studies, Harvard University. (2nd expanded ed.,
UCLA Asian Pacific Monograph Series, 2001.)
\item Fine, K. (2012). ``Guide to Ground.'' In Correia \& Schnieder
(eds.), 37--80.
\item Floridi, L. (2002). \emph{Sextus Empiricus: The Transmission and
Recovery of Pyrrhonism}. American Classical Studies 46. Oxford
University Press.
\item Fosl, P.S. (1998). ``Review of Norton \& Norton, \emph{The David
Hume Library}.'' \emph{ECSSS Newsletter} 11: 35--36.
\item Frede, M. (1983). ``Stoics and Skeptics on Clear and Distinct
Impressions.'' In M.~Burnyeat (ed.), \emph{The Skeptical Tradition},
65--93. University of California Press.
\item Garrett, D. (1997). \emph{Cognition and Commitment in Hume's
Philosophy}. Oxford University Press.
\item Goldman, A.I. (1979). ``What Is Justified Belief?'' In G.~Pappas
(ed.), \emph{Justification and Knowledge}, 1--23. D.~Reidel.
\item Goldman, A.I. (1986). \emph{Epistemology and Cognition}. Harvard
University Press.
\item Graham, A.C. (1978). \emph{Later Mohist Logic, Ethics and
Science}. Chinese University Press, Hong Kong.
\item Graham, A.C. (1989). \emph{Disputers of the Tao: Philosophical
Argument in Ancient China}. Open Court.
\item Hansen, C. (1983). \emph{Language and Logic in Ancient China}.
University of Michigan Press.
\item Hansen, C. (1992). \emph{A Daoist Theory of Chinese Thought: A
Philosophical Interpretation}. Oxford University Press.
\item Herbert of Cherbury. (1624). \emph{De Veritate}. Paris.
\item Hicks, R.D. (tr.) (1925). \emph{Diogenes Laertius: Lives of
Eminent Philosophers}, 2 vols. Loeb Classical Library.
\item Hume, D. (1739--40). \emph{A Treatise of Human Nature}. Norton \&
Norton (eds.), Oxford University Press, 2000.
\item Hume, D. (1748). \emph{An Enquiry Concerning Human Understanding}.
Beauchamp (ed.), Oxford University Press, 1999.
\item Hume, D. (1975). \emph{Enquiries}. Selby-Bigge \& Nidditch
(eds.), 3rd ed. Clarendon Press. (SBN pagination.)
\item Hunt, T.J. (1998). \emph{A Textual History of Cicero's Academici
Libri}. Brill.
\item Kjellberg, P. (1996). ``Sextus Empiricus, Zhuangzi, and Xunzi on
`Why Be Skeptical?'.'' In P.~Kjellberg \& P.J.~Ivanhoe (eds.),
\emph{Essays on Skepticism, Relativism, and Ethics in the Zhuangzi}.
SUNY Press.
\item Lagr\'ee, J. (1994). \emph{Juste Lipse et la restauration du
sto\"icisme}. Vrin.
\item Leibniz, G.W. (1714). \emph{Monadologie} \S32.
\item Lipsius, J. (1584). \emph{De Constantia}. Antwerp: Plantin.
\item Lipsius, J. (1604). \emph{Manuductionis ad Stoicam philosophiam
libri tres}; \emph{Physiologiae Stoicorum libri tres}. Antwerp:
Plantin-Moretus.
\item Locke, J. (1689). \emph{An Essay Concerning Human Understanding}
I.3.15--27. Nidditch (ed.), Clarendon Press, 1975.
\item Long, A.A. \& Sedley, D.N. (1987). \emph{The Hellenistic
Philosophers}, vol.~1. Cambridge University Press. Translation
``cognition'' for katalepsis at section 40; fist analogy at section
41B.
\item Millican, P. (2002). ``The Context, Aims, and Structure of Hume's
First Enquiry.'' In P.~Millican (ed.), \emph{Reading Hume on Human
Understanding}, 27--65. Clarendon Press.
\item Nawar, T. (2022). ``Clear and Distinct Perception in the Stoics,
Augustine, and William of Ockham.'' \emph{Aristotelian Society
Supplementary Volume} 96(1): 185--207.
\item Nidditch, P.H. (ed.) (1975). \emph{An Essay Concerning Human
Understanding}, by John Locke. Clarendon Press.
\item Norton, D.F. (1981). ``The Myth of `British Empiricism'.''
\emph{History of European Ideas} 1(4): 331--344.
\item Norton, D.F. \& Norton, M.J. (eds.) (1996). \emph{The David Hume
Library}. Edinburgh Bibliographical Society.
\item Popkin, R.H. (1951). ``David Hume: His Pyrrhonism and His
Critique of Pyrrhonism.'' \emph{Philosophical Quarterly} 1(5):
385--407.
\item Popkin, R.H. (1952). ``David Hume and the Pyrrhonian
Controversy.'' \emph{Review of Metaphysics} 6(1): 65--81. (Repr.\ in
\emph{The High Road to Pyrrhonism}, 133--147.)
\item Popkin, R.H. (1979). \emph{The History of Scepticism from Erasmus
to Spinoza}. Rev.\ ed. University of California Press.
\item Priest, G. (2010). ``The Logic of the Catu\c{s}ko\d{t}i.''
\emph{Comparative Philosophy} 1(2): 24--54.
\item Priest, G. (2018). \emph{The Fifth Corner of Four: An Essay
on Buddhist Metaphysics and the Catu\c{s}ko\d{t}i}. Oxford University Press.
\item Pruss, A.R. (2006). \emph{The Principle of Sufficient Reason: A
Reassessment}. Cambridge University Press.
\item Ro\v{s}ker, J. ``Epistemology in Chinese Philosophy.'' Stanford
Encyclopedia of Philosophy (accessed 2026).
\item Baltzly, D., Durand, M. \& Shogry, S. (2023). ``Stoicism.''
Stanford Encyclopedia of Philosophy (accessed 2026).
\item Bolyard, C. (2009; rev. 2025). ``Medieval Skepticism.'' Stanford
Encyclopedia of Philosophy (accessed 2026).
\item Papy, J. (2004; rev. 2019). ``Justus Lipsius.'' Stanford
Encyclopedia of Philosophy (accessed 2026).
\item Van Norden, B. ``Wang Yangming.'' Stanford Encyclopedia of
Philosophy (accessed 2026).
\item du Vair, G. (1594). \emph{De la constance et consolation \`es
calamit\'es publiques}. Paris.
\item Schmitt, C.B. (1972). \emph{Cicero Scepticus: A Study of the
Influence of the Academica in the Renaissance}. Martinus Nijhoff.
\item Schmitt, C.B. (1983). ``The Rediscovery of Ancient Skepticism in
Modern Times.'' In M.~Burnyeat (ed.), \emph{The Skeptical Tradition},
225--251. University of California Press.
\item Schnieder, B. (2011). ``A Logic for `Because'.'' \emph{Review of
Symbolic Logic} 4(3): 445--465. DOI:
10.1017/S1755020311000104.
\item Sextus Empiricus. (1562). \emph{Pyrrhoniae Hypotyposes}, tr.\
Henri Estienne. Geneva.
\item Sextus Empiricus. (1569). \emph{Adversus Mathematicos}, tr.\
Gentian Hervet. Paris: Martin Juvenis; Antwerp: Plantin.
\item Sextus Empiricus. (1621). \emph{Opera}, Greek ed.\ Pierre \&
Jacques Chou\"et. Geneva.
\item Sextus Empiricus. \emph{Adversus Mathematicos} VII;
\emph{Pyrrhoniae Hypotyposes} I. Bury (Loeb).
\item Tillemans, T.J.F. (1999). \emph{Scripture, Logic, Language:
Essays on Dharmak\={i}rti and His Tibetan Successors}. Wisdom
Publications.
\item Xunzi. \emph{Xunzi} 22, ``On the Rectification of Names''
(ctext.org; Knoblock tr., Stanford University Press, 1988--94).
\end{itemize}

\end{document}
